\chapter{电平、器件与数字抽象}

\begin{keyidea}
数字硬件不是从符号 0 和 1 开始，而是从可测量的电压区间、可控制的电流路径和明确的失效边界开始。本章建立从物理节点到可靠 bit 的第一层抽象。
\end{keyidea}

一块控制板上有一个启动按钮：按下时信号线接到电源，松开时信号线哪里都不接。写软件的人会以为松开就是 0，但那根悬空的线可能停在旧电压上，也可能被旁边信号线、人体或漏电流慢慢推走——接收端今天读到 0，明天读到 1。

贯穿案例是一块简化的安全启动控制板：

\begin{longtable}{L{0.22\linewidth}L{0.28\linewidth}L{0.4\linewidth}}
\toprule
信号 & 含义 & 需要检查的硬件问题 \\
\midrule
\code{start_button} & 人按下启动按钮 & 松开时是否悬空，按下时是否可靠为 1 \\\tablerowrule
\code{cover_closed} & 保护盖已经闭合 & 断线或接触不良时应回到安全默认值 \\\tablerowrule
\code{emergency_ok} & 急停没有被按下 & 信号名为正逻辑，但实物开关可能采用低有效接法 \\\tablerowrule
\code{temp_alarm} & 温度报警 & 任一报警都应阻止启动并点亮故障灯 \\\tablerowrule
\code{current_alarm} & 电流报警 & 报警输入不能依赖偶然电平或悬空节点 \\\tablerowrule
\code{allow_start} & 允许启动 & 必须由所有安全条件共同决定 \\\tablerowrule
\code{motor_enable_cmd} & 电机使能命令 & 只能由经过检查确认的组合结果产生 \\\tablerowrule
\code{fault_lamp} & 故障灯 & 需要汇总多个故障信号 \\
\bottomrule
\end{longtable}

\begin{keyidea}
数字硬件的第一层抽象是\hwkey{带余量的电平}，第一层实现是\hwkey{受控电流路径}。门、选择器、译码器都建立在这两层之上。
\end{keyidea}

\begin{pdfdiagram}
  \node[hw lane, minimum width=2.55cm, minimum height=1.05cm] (level) at (0,0) {电压区间\\可靠 0/1};
  \node[hw lane, minimum width=2.55cm, minimum height=1.05cm] (path) at (3.15,0) {电流路径\\上拉/下拉};
  \node[hw lane, minimum width=2.55cm, minimum height=1.05cm] (gate) at (6.3,0) {逻辑门\\真值表};
  \node[hw lane, minimum width=2.55cm, minimum height=1.05cm] (comb) at (9.45,0) {组合网络\\延迟/负载};
  \draw[hw bus] (level) -- (path);
  \draw[hw bus] (path) -- (gate);
  \draw[hw bus] (gate) -- (comb);
\end{pdfdiagram}
\diagramnote{从电平到组合逻辑的抽象链：后一层必须由前一层的物理事实支撑。}

\section{一、电平与噪声余量}

电路里说“一根线是 1”，完整含义不是“线上写着数字 1”，而是“这个节点相对于参考地的电压落在接收端承认的高电平区间内”。节点可以理解成多条金属连接汇合成的同一个电气位置；参考地是整块电路共同约定的 0V 参考点；电压是两个点之间的差值。数字电路里说某根信号线是 0V、3.3V 或 5V，实际都是在说它相对于地的电压。

电流说明电荷正在沿通路流动。若输出节点被一条通路接向电源，它可能被拉到高电平；若被一条通路接向地，它可能被拉到低电平。若两条通路都断开，节点不会自动变成 0，而是可能暂时停在旧电压附近，再被漏电流、干扰、后级输入或测量探针慢慢推走。若上拉和下拉两条强通路同时打开，电源到地之间形成冲突电流，电平不可靠，还会增加功耗和发热。

需要先区分导线和程序变量。变量没有被赋值时，语言可以规定默认值；导线没有被驱动时，只剩物理过程。它可能因为寄生电容保存旧电压，表现为保留上一次状态；也可能被旁边切换的信号线、人体接近、探针输入阻抗或后级漏电流轻微影响。调试数字硬件时，不能把这种偶然稳定误认为设计正确。

以按钮输入为例。假设按钮按下表示“请求发生”。一种不完整的接法是按钮一端接信号线，另一端接电源：

\begin{lstlisting}
button pressed  -> signal connected to VCC
button released -> signal connected to nothing
\end{lstlisting}

按下时，信号线被电源路径拉高，该部分行为符合预期。松开时，信号线并没有被明确拉低，它只是断开了。这个节点可能保留先前电荷，也可能被旁边导线耦合出一定电压，还可能被后级输入漏电流拖向某个不确定位置。于是接收端可能读到 1、0 或过渡电压。问题的根源不是程序错误，而是电路没有为节点规定明确归宿。

正确的第一层想法是：任何要被后级判断的输出节点，都必须在每一种输入情况下被明确地拉向高电平或低电平。按钮松开时需要一条弱下拉路径把信号线拉向地；按钮按下时，按钮路径把它拉向电源。弱下拉不是为了在按钮按下时和按钮争夺电源，而是为了在按钮松开时给节点一个默认归宿。它太弱，节点下降慢、抗干扰差；它太强，按钮按下时浪费电流，甚至让高电平不够高。

\begin{circuit}[x=1cm,y=1cm]
  % 左：按钮松开即悬空
  \node[hw label] at (0,2.6) {VCC};
  \draw[line width=0.55pt] (0,2.4) -- (0,1.95);
  \draw[fill=BookInk] (0,1.9) circle (0.05);
  \draw[fill=BookInk] (0,1.25) circle (0.05);
  \draw[line width=0.55pt] (-0.3,2.05) -- (0.3,1.75);
  \node[hw label] at (0.85,1.95) {按钮};
  \draw[line width=0.55pt] (0,1.25) -- (0,0.6);
  \draw[fill=BookInk] (0,0.6) circle (0.05);
  \draw[line width=0.55pt] (0,0.6) -- (1.5,0.6);
  \node[hw label] at (1.5,0.95) {到后级输入};
  \node[hw label, text=BookRed] at (0.35,-0.25) {松开时：悬空？};
  \node[hw label, text=BookRed, font=\small\bfseries] at (0.5,3.1) {错误接法};
  % 右：弱下拉给默认归宿
  \node[hw label] at (6,2.6) {VCC};
  \draw[line width=0.55pt] (6,2.4) -- (6,1.95);
  \draw[fill=BookInk] (6,1.9) circle (0.05);
  \draw[fill=BookInk] (6,1.25) circle (0.05);
  \draw[line width=0.55pt] (5.7,2.05) -- (6.3,1.75);
  \node[hw label] at (6.85,1.95) {按钮};
  \draw[line width=0.55pt] (6,1.25) -- (6,0.6);
  \draw[fill=BookInk] (6,0.6) circle (0.05);
  \draw[line width=0.55pt] (6,0.6) -- (7.5,0.6);
  \node[hw label] at (7.5,0.95) {到后级输入};
  \draw[line width=0.55pt] (6,0.6) -- (6,0.42);
  \draw[line width=0.55pt] (5.82,0.42) rectangle (6.18,-0.42);
  \draw[line width=0.55pt] (6,-0.42) -- (6,-0.55);
  \pic at (6,-0.55) {hw gnd};
  \node[hw label] at (7.15,0) {弱下拉电阻};
  \node[hw label, text=BookTeal, font=\small\bfseries] at (6.5,3.1) {正确接法};
\end{circuit}
\diagramnote{同一个按钮，两种结局。左图松开按钮后节点没有任何归宿，只能听天由命；右图多了一条弱下拉电阻：松开时它把节点拉向地（可靠 0），按下时按钮的强路径“赢过”弱下拉，把节点拉向电源（可靠 1）。}

\begin{longtable}{L{0.18\linewidth}L{0.31\linewidth}L{0.31\linewidth}L{0.12\linewidth}}
\toprule
场景 & 高电平路径 & 低电平路径 & 结论 \\
\midrule
按下 & 按钮把节点接向电源 & 弱下拉仍存在 & 可靠高 \\\tablerowrule
松开 & 高电平路径断开 & 弱下拉把节点接向地 & 可靠低 \\\tablerowrule
接触抖动 & 高电平路径快速断续 & 弱下拉持续存在 & 需要消抖 \\\tablerowrule
导线断裂 & 按钮路径永远断开 & 若下拉还在，则为低 & 有默认值 \\\tablerowrule
没有下拉 & 高电平路径断开 & 低电平路径也断开 & 悬空 \\
\bottomrule
\end{longtable}

把按钮放回安全启动控制板中。启动按钮松开时，系统不应因为线缆悬空而误认为有人按下；保护盖传感器断线时，系统不应默认保护盖闭合；报警线断开时，系统不应静默地继续允许运行。因此，“默认值”不是随意选择，而是安全策略的一部分。对人机输入而言，常见策略是让未动作、断线或未知状态对应不启动、不使能或报警一侧；对内部数据通路而言，策略可能不同，但仍必须明确。

\begin{longtable}{L{0.24\linewidth}L{0.28\linewidth}L{0.38\linewidth}}
\toprule
输入线 & 不完整设计的风险 & 更稳妥的默认方向 \\
\midrule
启动按钮 & 松开时悬空，可能误触发 & 默认不启动，按下才变为启动请求 \\\tablerowrule
保护盖检测 & 断线时可能被误读为闭合 & 默认未闭合，需要传感器明确证明闭合 \\\tablerowrule
急停链路 & 接触不良时可能丢失急停信号 & 默认不允许运行，需要链路明确证明急停正常 \\\tablerowrule
报警输入 & 断线时可能丢失报警 & 根据安全需求选择默认报警或要求上层诊断 \\
\bottomrule
\end{longtable}

电平也不能理解成两个数学点。假设一块小电路用 5V 供电。线上的真实电压可能是 0V、0.17V、1.3V、2.6V、4.8V，也可能在切换过程中从 0V 慢慢升到 5V。硬件不能要求每一级都精确读出这个连续数值。更可靠的做法是规定两个稳定区间：低到足够接近地电位时解释为 0，高到足够接近电源电位时解释为 1，中间一段不作为长期稳定状态。

\begin{longtable}{L{0.24\linewidth}L{0.28\linewidth}L{0.38\linewidth}}
\toprule
电压范围 & 逻辑解释 & 设计意义 \\
\midrule
接近地电位 & 逻辑 0 & 小噪声不会轻易把它推成 1 \\\tablerowrule
中间区域 & 不作为稳定状态 & 切换可以经过，但不能长期停留 \\\tablerowrule
接近电源电位 & 逻辑 1 & 小噪声不会轻易把它拉成 0 \\
\bottomrule
\end{longtable}

\begin{pdfdiagram}[x=1cm,y=1cm]
  \draw[BookRule,line width=0.5pt] (0,0) -- (10.8,0);
  \foreach \x/\label in {0/0V,2.0/$V_{IL}$,3.0/$V_{IH}$,5.0/$V_{CC}$} {
    \draw[BookMuted,line width=0.45pt] (2.0*\x,0.12) -- (2.0*\x,-0.12);
    \node[hw label] at (2.0*\x,-0.42) {\label};
  }
  \draw[BookTeal,line width=5pt] (0,0.45) -- (4.0,0.45);
  \draw[BookAmber,line width=5pt] (4.0,0.45) -- (6.0,0.45);
  \draw[BookAccent,line width=5pt] (6.0,0.45) -- (10.0,0.45);
  \node[BookTeal,font=\small\bfseries] at (2.0,0.95) {可靠 0};
  \node[BookAmber,font=\small\bfseries] at (5.0,0.95) {不确定区};
  \node[BookAccent,font=\small\bfseries] at (8.0,0.95) {可靠 1};
  \node[hw label, text width=3.1cm] (drv0) at (1.9,-1.25) {前级低电平必须低于后级低阈值};
  \node[hw label, text width=3.1cm] (drv1) at (8.1,-1.25) {前级高电平必须高于后级高阈值};
  \draw[hw return] (drv0.north) -- (2.0,-0.05);
  \draw[hw return] (drv1.north) -- (6.0,-0.05);
\end{pdfdiagram}
\diagramnote{电压不是两个数学点，而是带噪声余量的区间。中间区域允许切换经过，不允许作为稳定工作点。}

可以用四个电压指标描述两级电路之间的约定。$V_\mathrm{OH}$ 是输出端保证的高电平下限，$V_\mathrm{OL}$ 是输出端保证的低电平上限；$V_\mathrm{IH}$ 是输入端承认为高的最低电压，$V_\mathrm{IL}$ 是输入端承认为低的最高电压。若一个电路保证 $V_\mathrm{OH}$ 高于下一级需要的 $V_\mathrm{IH}$，中间差额就是高电平噪声余量；若 $V_\mathrm{OL}$ 低于下一级允许的 $V_\mathrm{IL}$，中间差额就是低电平噪声余量。

\begin{longtable}{L{0.18\linewidth}L{0.32\linewidth}L{0.4\linewidth}}
\toprule
符号 & 含义 & 读法 \\
\midrule
$V_\mathrm{OH}$ & 输出保证为高时的最差电压 & 前一级说：我给出的 1 至少高到这里 \\\tablerowrule
$V_\mathrm{IH}$ & 输入承认为高的最低电压 & 后一级说：至少给到这里，我才认作 1 \\\tablerowrule
$V_\mathrm{OL}$ & 输出保证为低时的最差电压 & 前一级说：我给出的 0 至多高到这里 \\\tablerowrule
$V_\mathrm{IL}$ & 输入承认为低的最高电压 & 后一级说：低到这里以下，我才认作 0 \\
\bottomrule
\end{longtable}

例如，某一级输出保证高电平至少为 2.4V，后一级输入只要求达到 2.0V 就承认为高，那么高电平噪声余量为 0.4V。若同一输出在最坏负载、温度和工艺条件下只能到 1.9V，逻辑表达式仍然可能写着“输出为 1”，但后一级不会可靠承认这个 1。低电平也同理：前级输出低电平的最坏值必须低于后一级允许的低电平上限，并留出余量。

\begin{longtable}{L{0.22\linewidth}L{0.22\linewidth}L{0.24\linewidth}L{0.22\linewidth}}
\toprule
项目 & 前级保证 & 后级要求 & 结论 \\
\midrule
高电平 & $V_\mathrm{OH} \ge 2.4V$ & $V_\mathrm{IH} \ge 2.0V$ & 余量 0.4V \\\tablerowrule
低电平 & $V_\mathrm{OL} \le 0.4V$ & $V_\mathrm{IL} \le 0.8V$ & 余量 0.4V \\\tablerowrule
失败情形 & 实测高电平仅 1.9V & 后级需要 2.0V & 不能可靠读作 1 \\
\bottomrule
\end{longtable}

把这个例子进一步写成定量算式，会更接近数据手册读法。假设某输出在最坏负载下保证 $V_\mathrm{OH}=3.0V$、$V_\mathrm{OL}=0.25V$，后级输入要求 $V_\mathrm{IH}=2.1V$、$V_\mathrm{IL}=0.7V$，则高电平噪声余量为 $3.0V-2.1V=0.9V$，低电平噪声余量为 $0.7V-0.25V=0.45V$。若板级串扰、地弹或线缆压降可能带来约 0.3V 扰动，这个连接仍有余量；若某条长线在最坏情况下可能叠加 0.6V 扰动，低电平余量已经不足，必须改变驱动、布线、阈值或输入调理方案。

\begin{longtable}{L{0.22\linewidth}L{0.24\linewidth}L{0.22\linewidth}L{0.22\linewidth}}
\toprule
检查项 & 算式 & 结果 & 解释 \\
\midrule
高电平余量 & $3.0V-2.1V$ & 0.9V & 前级高电平高于后级高阈值 \\\tablerowrule
低电平余量 & $0.7V-0.25V$ & 0.45V & 前级低电平低于后级低阈值 \\\tablerowrule
0.3V 干扰 & 余量仍为正 & 可接受 & 仍需用最坏条件确认 \\\tablerowrule
0.6V 干扰 & 低余量不足 & 不合格 & 可能把低电平推入不确定区 \\
\bottomrule
\end{longtable}

只要所有连接都满足这些最差条件，后续讨论 bit、真值表和指令编码才有可靠基础。若这些条件不满足，系统层面可能出现间歇性现象：同一块板有时能启动、有时不能；按键有时触发两次；某条总线在高温下出错；更换更长排线后原本稳定的电路开始误判。这些现象不一定来自逻辑表达式错误，也可能来自电平约定被破坏。

输出还必须能驱动下一级。一个节点如果只能在空载时看起来像 1，一接上后级就被拖低，它不是可靠数字输出。后级输入、板级走线、封装和探针都有寄生电容。把一个节点从 0 拉到 1，本质上是在给这个电容充电；把它从 1 拉到 0，本质上是在放电。驱动路径越弱，等效电阻越大；负载越多，等效电容越大；上升和下降越慢。

\begin{lstlisting}
larger load capacitance -> slower edge
weaker driver           -> slower edge
slower edge             -> less time margin before sampling
\end{lstlisting}

所以示波器上看到的真实信号不是理想直角。电压边沿有斜率，可能有过冲、振铃、平台和噪声。逻辑 0/1 的背后是一条随时间变化的电压曲线，接收端只是在某个时刻把它归类为 0 或 1。若这个时刻到来之前电压还没有进入可靠区间，即使最终稳定也不够。

因此，一个 bit 被称为可靠，至少需要同时满足四项条件。第一，有明确驱动路径，不能悬空、保留旧值或被干扰改变；第二，电平进入接收端承认的稳定区间并留有噪声余量，最坏输出也要满足后级阈值；第三，驱动端能够带动实际负载，不能让边沿被电容拖慢、电压被拖低；第四，在后级采样或使用之前，信号已经稳定足够长时间——最终值正确但采样时刻错误，同样是失败。

外部输入还需要调理。按钮和机械触点会抖动，线缆会引入噪声，传感器输出可能是开漏或开集结构，板外信号还可能与板内电源域不同。进入组合逻辑核心之前，外部输入应先变成带默认值、带驱动能力、带明确有效电平的内部信号。否则真值表虽然正确，输入源本身却可能不可靠。

\begin{longtable}{L{0.2\linewidth}L{0.34\linewidth}L{0.36\linewidth}}
\toprule
外部现象 & 进入组合逻辑前的处理 & 不处理的后果 \\
\midrule
按钮抖动 & 消抖或在后续同步边界过滤 & 一次按下可能产生多次请求 \\\tablerowrule
长线噪声 & 上拉/下拉、滤波、缓冲或屏蔽 & 误触发、边沿变慢、阈值附近摆动 \\\tablerowrule
开漏输出 & 外部上拉并限制总线负载 & 释放时没有可靠高电平 \\\tablerowrule
电源域不同 & 电平转换或隔离 & 后级阈值不匹配甚至损坏器件 \\
\bottomrule
\end{longtable}

板外输入还要有保护边界。按钮、传感器、线缆和连接器会把静电放电、误接电压、浪涌、长线耦合和地电位差带到芯片引脚附近。若把这些信号直接接入 CMOS 输入，逻辑上也许只是“读一个 0/1”，电气上却可能让输入保护二极管承受过大电流，或者让引脚在绝对最大额定值之外工作。保护电路的职责不是改变业务含义，而是在进入逻辑核心之前限制\hwrisk{电压、电流、带宽和边沿形态}。

\begin{longtable}{L{0.2\linewidth}L{0.34\linewidth}L{0.36\linewidth}}
\toprule
保护措施 & 主要作用 & 设计时必须说明 \\
\midrule
串联电阻 & 限制异常电压下的输入电流，减小尖峰冲击 & 电阻不能大到让边沿过慢或阈值附近停留过久 \\\tablerowrule
RC 滤波 & 抑制低速按钮抖动和高频噪声 & 时间常数要匹配输入速度，不能滤除有效脉冲 \\\tablerowrule
TVS 或 ESD 保护 & 为静电和瞬态浪涌提供旁路 & 器件选型要匹配工作电压、钳位电压和能量等级 \\\tablerowrule
输入钳位与限流 & 让引脚异常电压不直接伤害芯片内部结构 & 不能依赖片内保护长期承受外部错误连接 \\\tablerowrule
缓冲或施密特输入 & 恢复边沿、提供滞回、隔离前级负载 & 输出电平、传播延迟和驱动能力仍需逐项确认 \\
\bottomrule
\end{longtable}

因此，板外按钮或传感器不应被写成“直接连到门输入”。更严谨的写法是：连接器输入先经过默认上拉/下拉、限流、保护、滤波或施密特边界，再形成内部语义信号。这样即使后续表达式仍是 \code{allow_start = start_request AND cover_closed AND emergency_ok AND no_fault}，读者也能看出 \code{start_request} 已经不是裸露引脚，而是经过电气边界整理后的可靠 bit。

\topic{从外部引脚到内部 bit 的五层边界}
外部世界进入芯片内部之前，至少要经过五层边界。每一层都回答一个不同问题：连接器和线缆负责把物理事件带到板上；保护与限流负责让异常能量不直接伤害芯片；阈值整形负责把连续电压变成稳定 0/1；语义归一化负责把低有效、反相、断线默认和消抖结果整理成内部名字；后级逻辑只应依赖最后得到的内部语义信号。

\begin{pdfdiagram}
  \node[hw block, text width=2.1cm, minimum height=1.15cm] (raw) at (0,0) {外部物理量\\按钮/线缆};
  \node[hw block, text width=2.1cm, minimum height=1.15cm] (protect) at (2.65,0) {保护边界\\限流/ESD};
  \node[hw block, text width=2.1cm, minimum height=1.15cm] (shape) at (5.3,0) {阈值整形\\滤波/施密特};
  \node[hw block, text width=2.1cm, minimum height=1.15cm] (sem) at (7.95,0) {语义归一\\极性/消抖};
  \node[hw block, text width=2.1cm, minimum height=1.15cm] (logic) at (10.6,0) {组合逻辑\\真值表};
  \draw[hw bus] (raw) -- (protect);
  \draw[hw bus] (protect) -- (shape);
  \draw[hw bus] (shape) -- (sem);
  \draw[hw bus] (sem) -- (logic);
  \node[hw label, text width=2.35cm] at (0,-1.15) {可抖动、可断线、可带噪声};
  \node[hw label, text width=2.35cm] at (2.65,-1.15) {限制电压、电流和瞬态能量};
  \node[hw label, text width=2.35cm] at (5.3,-1.15) {把慢边沿和噪声推出不确定区};
  \node[hw label, text width=2.35cm] at (7.95,-1.15) {得到正逻辑内部事实};
  \node[hw label, text width=2.35cm] at (10.6,-1.15) {只处理可靠 bit};
\end{pdfdiagram}
\diagramnote{外部输入进入组合逻辑前的五层边界。图中的箭头不是简单连线，而是逐层收窄不确定性的过程。}

这张图能避免两个常见错误。第一，不能把 \code{start_button_raw} 直接代入真值表，因为它还包含抖动、噪声、线缆和保护问题。第二，不能把保护电路当作逻辑表达式的一部分，因为保护电路的职责是限制物理风险，而不是决定业务语义。一个严谨的规格应分别写出 \code{start_button_raw}、\code{start_button_pin}、\code{start_button_clean} 和 \code{start_request}，并说明每个名字跨越了哪一层边界。

上拉电阻和 RC 滤波可以用一个数量级例题理解。若开漏信号采用 4.7k$\Omega$ 上拉到 3.3V，低电平被拉下时的电流约为 $3.3V/4.7k\Omega \approx 0.7mA$，后级通常较容易承受；若同一节点总电容约 100pF，则时间常数约为 $4.7k\Omega \times 100pF = 0.47\mu s$，低速按钮或报警线通常可以接受。若把上拉改成 100k$\Omega$，同样 100pF 下时间常数变成 10$\mu s$，边沿会明显变慢；若后级很快采样，或者输入没有施密特滞回，信号就可能在阈值附近停留太久。

\begin{longtable}{L{0.2\linewidth}L{0.28\linewidth}L{0.42\linewidth}}
\toprule
参数选择 & 数量级结果 & 设计含义 \\
\midrule
4.7k$\Omega$ 上拉到 3.3V & 拉低电流约 0.7mA & 要确认开漏器件能可靠吸收该电流 \\\tablerowrule
4.7k$\Omega$ 与 100pF & 时间常数约 0.47$\mu s$ & 适合低速输入，仍需看采样时刻 \\\tablerowrule
100k$\Omega$ 与 100pF & 时间常数约 10$\mu s$ & 边沿显著变慢，抗噪声能力下降 \\\tablerowrule
RC 消抖 10ms & 能过滤机械抖动 & 不适合需要微秒级响应的输入 \\
\bottomrule
\end{longtable}

这些数值不是固定推荐值，而是提醒读者把“能不能读到 1”拆成电流、边沿和时间三个问题。上拉过强会增加拉低电流，上拉过弱会让释放后的高电平上升太慢；RC 过小无法过滤抖动，RC 过大又会滤除有效边沿。真正设计时应以器件数据手册、输入速度、阈值、功耗和安全响应时间共同决定。

\topic{上拉/下拉电阻的阻值权衡}
“上拉取 10k$\Omega$”常被当作口诀背下来，但阻值其实是一笔可以算清的账。先看电阻太大的一端。后级输入并非完全不取电流，CMOS 输入存在微安量级以下的漏电流，而这股电流必须流过上拉电阻，在电阻两端形成压降 $V=IR$，于是高电平被拉低。以 5V 系统、10k$\Omega$ 上拉、输入漏电流 1$\mu A$ 为例：压降为 $1\mu A \times 10k\Omega = 10mV$，实测高电平约 4.99V，几乎无损。若为了“省电”把上拉换成 1M$\Omega$，同样 1$\mu A$ 就产生 1V 压降，高电平只剩 4.0V；若高温下漏电升到 10$\mu A$，压降将达到 10V 量级——节点根本到不了高电平。阻值越大，默认电平越怕漏电，噪声余量越薄。

再看电阻太小的一端：按键或开漏器件把节点拉低时，上拉电阻直接从电源向地通电。若上拉只有 100$\Omega$，按键按下期间持续流过 $5V/100\Omega = 50mA$，电阻上的功耗为 $I^2R = 0.05^2 \times 100 = 0.25W$，已超过常见 1/8W 贴片电阻的额定功率，按着不放就会发烫，拉低它的器件也必须灌得进这 50mA。换成 10k$\Omega$，同样动作只消耗 0.5mA，普通逻辑输出和开漏器件都能承受。

但阻值也不能无限大。默认电平靠电阻“拽住”，阻值越大节点越“轻”，同样的耦合电荷产生的电压扰动越大，抗干扰越差；同时电阻与节点电容构成 RC 充电回路，100k$\Omega$ 配 50pF 走线电容的时间常数已达 $100k\Omega \times 50pF = 5\mu s$。综合起来，阻值由灌/拉电流能力、噪声余量、功耗三者权衡，常用范围在 1k$\Omega$ 到 100k$\Omega$ 之间：需要强抗干扰或较快边沿时偏向下端，电池供电、信号很慢时偏向上端。

\begin{longtable}{L{0.24\linewidth}L{0.24\linewidth}L{0.16\linewidth}L{0.26\linewidth}}
\toprule
情形 & 算式 & 结果 & 说明 \\
\midrule
10k$\Omega$ 上拉，漏电 1$\mu A$ & $1\mu A \times 10k\Omega$ & 压降 10mV & 5V 高电平约 4.99V，裕量几乎无损 \\\tablerowrule
1M$\Omega$ 上拉，漏电 1$\mu A$ & $1\mu A \times 1M\Omega$ & 压降 1V & 高电平只剩 4.0V，高温下更差 \\\tablerowrule
100$\Omega$ 上拉，按键接地 & $5V/100\Omega$，$I^2R$ & 50mA、0.25W & 超过 1/8W 电阻额定，按着不放会发烫 \\\tablerowrule
10k$\Omega$ 上拉，按键接地 & $5V/10k\Omega$ & 0.5mA & 功耗可忽略，开漏器件也灌得起 \\\tablerowrule
100k$\Omega$ 上拉，节点 50pF & $100k\Omega \times 50pF$ & $\tau = 5\mu s$ & 边沿明显变慢，抗扰动能力下降 \\
\bottomrule
\end{longtable}

“先看数据手册再定阻值”的原因在于：输入漏电、开漏器件的灌电流能力和节点电容分别决定阻值的上限和下限，口诀只是在常见参数下的折中。

由此可见，书写硬件规格时应区分“原始外部信号”和“内部逻辑信号”。原始外部信号属于板外世界，可能带有抖动、噪声、线缆压降和接插件故障；内部逻辑信号属于组合逻辑世界，必须已经具有明确的 0/1 含义。若把二者混成一个名字，读者会误以为真值表直接约束了外部物理现象。

\begin{longtable}{L{0.24\linewidth}L{0.34\linewidth}L{0.32\linewidth}}
\toprule
信号层次 & 典型名称 & 设计要求 \\
\midrule
外部触点 & \code{start_button_raw} & 允许抖动，但必须有默认路径和保护边界 \\\tablerowrule
输入调理后 & \code{start_button_clean} & 在逻辑核心看来只能是稳定 0 或稳定 1 \\\tablerowrule
语义信号 & \code{start_request} & 明确表示“本周期有启动请求”这一逻辑事实 \\\tablerowrule
安全汇总 & \code{allow_start} & 由全部安全条件共同决定，不能依赖原始触点 \\
\bottomrule
\end{longtable}

施密特触发输入是常见的边界处理方式。普通输入通常只有一个判断阈值附近的过渡区；当信号边沿很慢或带噪声时，电压可能在阈值附近反复穿越，导致接收端多次翻转。施密特触发器使用两个不同阈值：上升时必须超过较高阈值才承认为 1，下降时必须低于较低阈值才承认为 0。两个阈值之间的间隔称为滞回。滞回不是改变逻辑功能，而是让输入在噪声环境中更稳定地跨过边界。

\begin{waveform}[x=1cm,y=1cm]
  \draw[hw bus] (0,0) -- (6.4,0);
  \draw[hw bus] (0,0) -- (0,3.3);
  \node[hw label] at (6.4,-0.35) {输入电压};
  \node[hw label, rotate=90] at (-0.55,1.6) {输出};
  \draw[BookMuted, dashed, line width=0.4pt] (2.2,0) -- (2.2,2.6);
  \draw[BookMuted, dashed, line width=0.4pt] (3.8,0) -- (3.8,2.6);
  \node[hw label] at (2.2,-0.35) {$V_{T-}$};
  \node[hw label] at (3.8,-0.35) {$V_{T+}$};
  \node[hw label] at (-0.5,2.5) {高};
  \node[hw label] at (-0.5,0.5) {低};
  \draw[BookAccent, line width=1.1pt] (0,0.55) -- (3.8,0.55) -- (3.8,2.55) -- (5.9,2.55);
  \draw[BookAccent, line width=1.1pt, -{Latex[length=2mm]}] (1.7,0.55) -- (2.3,0.55);
  \draw[BookAccent, line width=1.1pt, -{Latex[length=2mm]}] (3.8,1.35) -- (3.8,1.95);
  \draw[BookTeal, line width=1.1pt] (5.9,2.45) -- (2.2,2.45) -- (2.2,0.45) -- (0,0.45);
  \draw[BookTeal, line width=1.1pt, -{Latex[length=2mm]}] (4.3,2.45) -- (3.7,2.45);
  \draw[BookTeal, line width=1.1pt, -{Latex[length=2mm]}] (2.2,1.75) -- (2.2,1.15);
  \node[hw label, text=BookAccent, text width=2.6cm] at (5.0,1.15) {上升：越过 $V_{T+}$ 才承认 1};
  \node[hw label, text=BookTeal, text width=2.6cm] at (1.0,3.05) {下降：低于 $V_{T-}$ 才承认 0};
\end{waveform}
\diagramnote{施密特输入的传输特性曲线。普通输入只有一个模糊阈值，噪声会让输出在阈值附近反复横跳；施密特把“上升承认”和“下降承认”拉开成两个阈值，中间这段滞回区间就是噪声的容身之地——信号必须真正越过 $V_{T+}$ 才算 1，真正跌回 $V_{T-}$ 以下才算 0。}

\begin{longtable}{L{0.22\linewidth}L{0.34\linewidth}L{0.34\linewidth}}
\toprule
输入处理 & 适用情形 & 对可靠 bit 的贡献 \\
\midrule
上拉或下拉 & 按钮、开漏输出、断线默认值 & 给悬空节点明确归宿 \\\tablerowrule
限流和保护 & 板外连接、误插、瞬态干扰 & 避免异常电压直接进入逻辑核心 \\\tablerowrule
滤波或消抖 & 机械触点、低速传感器 & 把多次物理跳变收敛为一次逻辑事件 \\\tablerowrule
施密特触发 & 慢边沿、噪声叠加的输入 & 用滞回减少阈值附近的反复翻转 \\\tablerowrule
电平转换 & 1.8V、3.3V、5V 等混合系统 & 让输出约定匹配输入约定 \\
\bottomrule
\end{longtable}

按钮消抖尤其适合作为入门反例。人只按下一次，触点却可能在几毫秒内多次接通、断开。如果后级把每一次跳变都看成一次请求，一个启动按钮就可能产生多个启动脉冲。消抖电路或后续时序逻辑的职责，是在确认输入持续稳定后才更新内部请求信号。

\begin{longtable}{L{0.18\linewidth}L{0.22\linewidth}L{0.24\linewidth}L{0.26\linewidth}}
\toprule
时间段 & 原始触点 & 内部请求 & 解释 \\
\midrule
按下前 & 稳定断开 & 0 & 默认下拉给出明确低电平 \\\tablerowrule
刚按下 & 断续跳变 & 保持 0 & 尚未接受为一次稳定请求 \\\tablerowrule
稳定按下 & 持续接通 & 1 & 已满足消抖条件 \\\tablerowrule
刚松开 & 断续跳变 & 保持 1 或等待释放确认 & 不把抖动解释为多次按键 \\\tablerowrule
稳定松开 & 持续断开 & 0 & 回到默认状态 \\
\bottomrule
\end{longtable}

\begin{waveform}[x=1cm,y=1cm]
  \draw[hw bus] (0,0) -- (12.4,0);
  \node[hw label] at (12.4,-0.35) {时间};
  \node[hw label, anchor=east] at (-0.2,2.4) {原始触点};
  \node[hw label, anchor=east] at (-0.2,0.4) {消抖后的内部请求};
  \draw[BookAccent, line width=1.0pt] (0,2.0) -- (2,2.0) -- (2,2.8) -- (2.25,2.8) -- (2.25,2.0) -- (2.5,2.0) -- (2.5,2.8) -- (2.75,2.8) -- (2.75,2.0) -- (3,2.0) -- (3,2.8) -- (8,2.8) -- (8,2.0) -- (8.25,2.0) -- (8.25,2.8) -- (8.5,2.8) -- (8.5,2.0) -- (12,2.0);
  \draw[BookTeal, line width=1.1pt] (0,0) -- (4,0) -- (4,0.8) -- (9.4,0.8) -- (9.4,0) -- (12,0);
  \draw[BookMuted, dashed, line width=0.4pt] (2,-0.3) -- (2,3.1);
  \draw[BookMuted, dashed, line width=0.4pt] (3.4,-0.3) -- (3.4,3.1);
  \node[hw label] at (2.7,3.35) {抖动区};
  \draw[BookMuted, dashed, line width=0.4pt] (8,-0.3) -- (8,3.1);
  \draw[BookMuted, dashed, line width=0.4pt] (8.9,-0.3) -- (8.9,3.1);
  \node[hw label] at (8.45,3.35) {抖动区};
  \draw[hw return] (4,1.05) -- node[hw label, above]{确认稳定后才置位} (6.6,1.05);
  \node[hw label, text=BookTeal] at (6.7,0.4) {一次按下只产生一次请求};
\end{waveform}
\diagramnote{人只按了一次，触点却在几毫秒内多次通断（抖动区）。消抖的逻辑是“连续稳定一段时间才算数”：内部请求比真实按下晚一点点，代价可忽略，换来的是它只跳变一次。}

这组处理说明一个原则：逻辑核心应当面对已经被整理过的事实，而不是直接面对物理噪声。安全启动控制板中，\code{allow_start} 不应直接依赖 \code{start_button_raw}，而应依赖 \code{start_request}；故障汇总不应直接依赖线缆上的未定义电平，而应依赖经过默认值、阈值和驱动能力检查后的报警信号。

可靠 bit 还必须考虑上电和复位。电源从 0V 上升到额定电压需要时间，芯片内部阈值、外部上拉下拉、时钟源和输入调理电路不一定同时进入正常状态。在这段时间里，若某个输出直接控制电机、继电器或写使能，就不能假设它已经是安全值。正式硬件设计通常要求复位链路或电源监控电路在电源稳定之前保持关键输出禁止。

\begin{longtable}{L{0.18\linewidth}L{0.34\linewidth}L{0.38\linewidth}}
\toprule
阶段 & 硬件状态 & 对关键输出的要求 \\
\midrule
断电 & 节点可能通过漏电或外部连接获得微弱电压 & 不应驱动执行机构 \\\tablerowrule
电源爬升 & 阈值、时钟和输入调理尚未完全有效 & 复位或禁止链路保持有效 \\\tablerowrule
复位保持 & 内部状态被强制到约定默认值 & \code{motor_enable_cmd}=0，写使能关闭 \\\tablerowrule
复位释放 & 输入和时钟达到有效条件 & 组合逻辑开始按规格产生输出 \\\tablerowrule
正常运行 & 所有电平约定和时序约定成立 & 后级只使用已经稳定的信号 \\
\bottomrule
\end{longtable}

这也是“默认安全”的电气含义。它不是在文档里写一句“默认不启动”，而是要能指出：电源未稳时谁拉住使能，复位未释放时谁关闭输出，输入断线时谁把条件解释为不满足，组合路径未稳定时谁阻止后级采样。若这些问题回答不出来，系统可能在最危险的上电瞬间出现短促误动作。

最后，可靠 bit 还要能被实际测量确认。万用表可以检查静态电压，却看不见窄脉冲和抖动；示波器能观察边沿、过冲、振铃和噪声；逻辑分析仪适合记录长时间数字事件，但它同样依据自己的阈值把模拟电压解释成 0/1。测量工具不是抽象之外的附属品，而是把电平约定落实到实物上的手段。

\begin{longtable}{L{0.2\linewidth}L{0.34\linewidth}L{0.36\linewidth}}
\toprule
工具或手段 & 能说明什么 & 不能单独证明什么 \\
\midrule
原理图审查 & 是否有默认路径、保护和驱动边界 & 实物焊接和噪声环境是否合格 \\\tablerowrule
万用表 & 静态高低电平是否大致正确 & 边沿速度、毛刺和短时抖动 \\\tablerowrule
示波器 & 波形、阈值穿越、过冲和振铃 & 长时间低概率事件一定不会发生 \\\tablerowrule
逻辑分析仪 & 数字事件顺序和持续时间 & 模拟电平余量是否足够 \\\tablerowrule
数据手册参数 & 最坏阈值、延迟、负载能力 & 板级实现是否满足这些边界 \\
\bottomrule
\end{longtable}

同一根按钮输入线，可以用三类工具得到互补的测量结果。假设 \code{start_button_raw} 采用上拉输入，按下时接地，内部逻辑再把它反相为 \code{start_request}。万用表能确认松开和按下时的静态电压方向；示波器能观察按下瞬间的触点抖动、边沿速度和阈值穿越；逻辑分析仪能记录内部请求是否只产生一次事件。三者关注的问题不同，不能互相替代。

\begin{longtable}{L{0.18\linewidth}L{0.33\linewidth}L{0.39\linewidth}}
\toprule
测量手段 & 在按钮案例中应看到什么 & 若结果异常应怀疑什么 \\
\midrule
万用表 & 松开为稳定高电平，按下为稳定低电平 & 上拉/下拉错误、接线反向、触点失效 \\\tablerowrule
示波器 & 按下和松开有有限抖动，最终越过阈值并稳定 & 抖动过长、边沿太慢、噪声跨越阈值 \\\tablerowrule
逻辑分析仪 & 消抖后只出现一次 \code{start_request} & 消抖不足、阈值设置不当、采样边界错误 \\\tablerowrule
数据手册 & 输入阈值、滞回、电流和最大额定值满足设计 & 电平约定只在实验样品上偶然成立 \\
\bottomrule
\end{longtable}

\topic{万用表与示波器各自能看见什么}
万用表直流电压档测出的是一段时间内的平均值。静态节点上，松开约 5V、按下约 0V，读数直接可信；但一个 0V/5V、占空比 50\% 的时钟信号，直流档读数约为 2.5V——既不是 0 也不是 1，无法据此判断波形好坏。窄毛刺、触点几毫秒的抖动、一次几纳秒的过冲，在平均值里几乎不留痕迹。所以万用表适合回答“默认电平对不对、供电对不对”，不适合回答“信号干不干净”。

示波器把电压随时间的曲线显示出来，才能看到边沿时间、过冲、振铃和阈值穿越次数。读边沿时要看 10\% 到 90\% 电平之间的行程时间，而不是只看最终稳定值；按钮消抖是否成功、复位释放是否单调、开漏上拉是否太慢，都要在波形上确认。两类工具回答的是不同问题，不能互相替代。

还要记住：测量仪器接入后本身就是被测电路的一部分。万用表和示波器的输入电阻通常约为 10M$\Omega$，去测一个靠 1M$\Omega$ 上拉维持的节点，相当于给节点并联了一个 10M$\Omega$ 电阻，分压后读数约为 $5V \times 10M\Omega/(1M\Omega+10M\Omega) \approx 4.55V$——比真实值低了约 0.45V，可能把一个勉强合格的高电平误判成故障。示波器探头还带有电容：$\times$10 探头约 10 到 15pF，$\times$1 探头可达上百 pF；把 15pF 加到 100k$\Omega$ 上拉的节点上，时间常数就多了 $100k\Omega \times 15pF = 1.5\mu s$，边沿会被明显拖慢。测量会改变被测对象，在高阻、高速节点上尤其明显。

\begin{longtable}{L{0.2\linewidth}L{0.36\linewidth}L{0.34\linewidth}}
\toprule
工具 & 能读出的内容 & 读不出或会引入的问题 \\
\midrule
万用表直流档 & 静态高低电平、供电电压、慢变化平均值 & 看不到毛刺、抖动和边沿，快速波形只剩平均值 \\\tablerowrule
示波器 & 边沿时间、过冲、振铃、阈值穿越、瞬态过程 & 探头电容与输入电阻会改变高阻节点的读数 \\\tablerowrule
测量行为本身 & 10M$\Omega$ 级输入电阻、pF 级探头电容 & 高阻节点被分压，边沿被附加电容拖慢 \\
\bottomrule
\end{longtable}

因此读数之前要先问两个问题：仪器接入是否改变了这个节点，以及读数代表的是哪一段时间内的电压。

还要补上板级物理边界。数字信号不是只沿着“信号线”单独传播；每一次电流变化都需要回流路径，供电和地也有电阻、电感和寄生电容。若多个输出同时翻转，瞬时电流会让地线或电源线上出现小幅电压变化，表现为地弹、供电跌落或参考点移动。若两根长线靠得太近，快速边沿还可能通过寄生电容或电感耦合到邻线，表现为串扰。它们不会改变布尔表达式，却会改变接收端看到的真实电压。

\begin{longtable}{L{0.2\linewidth}L{0.34\linewidth}L{0.36\linewidth}}
\toprule
板级现象 & 物理来源 & 对可靠 bit 的影响 \\
\midrule
回流路径过长 & 信号电流的返回路径绕远或不连续 & 边沿振铃、辐射增加、邻线耦合增强 \\\tablerowrule
地弹 & 多个输出同时切换，地线上产生瞬时压降 & 接收端参考地移动，噪声余量被压缩 \\\tablerowrule
供电跌落 & 瞬时电流超过局部供电支撑能力 & 输出驱动变弱，边沿变慢或阈值判断不稳 \\\tablerowrule
串扰 & 相邻走线之间存在寄生电容和互感 & 静止信号上出现错误脉冲或阈值附近扰动 \\\tablerowrule
长线振铃 & 走线具有传输线特性且终端不匹配 & 接收端可能多次穿越阈值 \\\tablerowrule
去耦不足 & 芯片附近缺少局部电荷支撑 & 切换时电源/地噪声增大 \\
\bottomrule
\end{longtable}

因此，电平约定至少隐含三类物理条件。第一，信号源和接收端必须共享可接受的参考地，或者通过隔离/差分/电平转换重新定义边界。第二，供电必须能在切换瞬间提供局部电流，去耦电容的职责就是在芯片附近提供短时间电荷支撑。第三，长线和外部连接器不能只按“理想导线”理解，必要时要考虑终端、屏蔽、滤波、缓冲和同步采样。“数字硬件基础”部分不展开 PCB 设计，但必须让读者知道：真值表正确并不自动保证板级物理边界正确。

\topic{输入模型和输出模型：引脚不是理想变量}
为了把原理讲清楚，可以把每个数字引脚暂时看成一个等效模型。输入端不是完全不取电流的数学变量，它有输入漏电、输入电容、阈值和绝对最大额定值；输出端也不是理想电压源，它有等效导通电阻、最大源电流/灌电流、传播延迟和短路风险。把这些模型放进脑中，很多“逻辑式没错但板子不稳”的问题就能被提前解释。

\begin{longtable}{L{0.18\linewidth}L{0.36\linewidth}L{0.36\linewidth}}
\toprule
引脚模型 & 关键参数 & 影响的问题 \\
\midrule
输入漏电 & 输入端即使不主动驱动，也可能有微小电流 & 上拉/下拉过弱时，默认电平可能被漏电和噪声改变 \\\tablerowrule
输入电容 & 输入门电容、封装和走线电容共同形成负载 & 边沿速度、RC 时间常数和阈值穿越时间 \\\tablerowrule
输入阈值 & $V_\mathrm{IH}$、$V_\mathrm{IL}$、滞回和电源域 & 电平余量、慢边沿误触发和电平转换需求 \\\tablerowrule
输出电阻 & 上拉/下拉晶体管导通后仍有等效电阻 & 负载越重，输出电平越偏离理想电源或地 \\\tablerowrule
输出电流 & 源电流、灌电流和总封装电流有限 & 扇出、LED/继电器驱动、总线冲突和发热 \\\tablerowrule
输出延迟 & 输入变化到输出越过阈值需要时间 & 组合路径最长延迟、毛刺宽度和采样窗口 \\
\bottomrule
\end{longtable}

因此，扇出不是“一个输出能画多少根线”的绘图问题，而是负载、电流和延迟问题。一个输出同时驱动十个 CMOS 输入，静态电流也许很小，但输入电容会叠加，边沿会变慢；同一个输出若还要驱动 LED、长线或继电器线圈，就已经超出普通逻辑门职责，通常需要缓冲器、晶体管、MOSFET、驱动芯片或隔离器。工程文档应明确区分“逻辑输出”和“功率驱动输出”。

\begin{longtable}{L{0.22\linewidth}L{0.34\linewidth}L{0.34\linewidth}}
\toprule
连接方式 & 不能只看逻辑式的原因 & 更稳妥的检查问题 \\
\midrule
一个门驱动多个输入 & 负载电容叠加，边沿变慢 & 最坏负载下是否按时越过阈值 \\\tablerowrule
逻辑门直接点亮 LED & LED 电流可能超过输出额定值 & 是否需要限流、电流预算和专用驱动 \\\tablerowrule
普通输出接长线 & 长线可能振铃、串扰或引入地电位差 & 是否需要缓冲、终端、滤波或差分传输 \\\tablerowrule
两个输出并接 & 一高一低时会形成强冲突电流 & 是否改用开漏、三态仲裁或总线协议 \\\tablerowrule
跨电源域连接 & 阈值和最大额定值可能不匹配 & 是否需要电平转换或隔离 \\
\bottomrule
\end{longtable}

读数据手册时，也应把这些模型翻译成检查清单：输入端看阈值、漏电、电容和最大额定值；输出端看高低电平保证、驱动电流、传播延迟、上升下降时间和总功耗限制。只要这张清单没有填完，布尔式就还不能直接等同于可工作的硬件。

\section{二、从器件到受控电流路径}

二极管可以先理解成方向性很强的器件。在合适偏置下，它比较容易让电流沿一个方向通过；反向偏置时，它基本阻止电流通过。这个模型不需要一开始展开半导体物理，却已经足够说明硬件设计的一条基本路线：利用器件物理特性，制造可预测的电流路径。

二极管的“单向”不是理想数学开关。正向导通通常需要一定压差，导通后还有压降；反向并非绝对没有电流，而是漏电很小；电压过高时还可能击穿。初读数字硬件时可以先用理想模型建立方向感，但不能忘记真实器件会消耗电压余量、引入延迟和功耗。

把数量级写下来，后面的讨论才有抓手。常见硅二极管的正向压降约为 0.6V 到 0.7V，锗二极管约为 0.2V 到 0.3V，肖特基二极管约为 0.2V 到 0.4V；反向漏电流通常在微安量级以下，普通小功率管的反向耐压从几十伏到几百伏不等。本章估算一律取硅管 0.7V：一路 5V 信号经过一只二极管后，高电平只剩约 4.3V；再串一级，就只剩约 3.6V。若后级输入高阈值是 2.0V，单看一级还剩 2.3V 余量，串到第三级时余量只剩 0.9V——这就是“级联多了电平余量越来越紧”的定量含义。

\begin{longtable}{L{0.24\linewidth}L{0.3\linewidth}L{0.36\linewidth}}
\toprule
级数 & 输出高电平（输入 5V） & 距 2.0V 阈值的余量 \\
\midrule
0 级（直连） & 5.0V & 3.0V \\\tablerowrule
1 级二极管 & $5.0-0.7=4.3$V & 2.3V \\\tablerowrule
2 级二极管 & $4.3-0.7=3.6$V & 1.6V \\\tablerowrule
3 级二极管 & $3.6-0.7=2.9$V & 0.9V \\
\bottomrule
\end{longtable}

\begin{longtable}{L{0.2\linewidth}L{0.34\linewidth}L{0.36\linewidth}}
\toprule
二极管事实 & 对数字电路的影响 & 读书时要追问 \\
\midrule
有方向性 & 可以限制电流路径 & 这条路径在何种输入下导通 \\\tablerowrule
有正向压降 & 输出电平会损失一部分余量 & 级联后高电平还够不够高 \\\tablerowrule
有漏电和极限 & 不能把反向阻断看成无限好 & 极端温度、电压下是否仍可靠 \\\tablerowrule
不能主动放大 & 不能恢复变弱的电平 & 下一级是否需要缓冲或开关恢复 \\
\bottomrule
\end{longtable}

\topic{整流、钳位与续流：二极管的三个工程角色}
在把二极管用进逻辑之前，先看它在电源与保护中的三个更常见的工程角色。这三种用法都只依赖“单向导通”这一条性质，解决的却是三类完全不同的事故。

第一个角色是整流：把交变或极性不定的电压整理成单向电流。接法上，可以单管串联做半波整流，也可以用四只二极管组成桥式全波整流，再配合电容和稳压得到直流。电源入口串一只二极管同时也是最低成本的防反接：电源装反时它不导通，电路只是不工作，而不是烧毁。失效后果是直接的——整流二极管接反或击穿短路后，反向电压或交流分量直接加到直流母线上，后级电解电容和稳压器可能损坏。

第二个角色是钳位保护：把引脚电压限制在电源轨附近。接法是用两只二极管从信号引脚分别“反并”到电源和地：引脚电压高于 $V_{CC}+0.7V$ 时上管导通，把电流引入电源轨；低于 $-0.7V$ 时下管导通。引脚因此被钳在电源轨上下各一个二极管压降的范围之内。芯片引脚内部的 ESD 保护正是这种结构。没有钳位时，人体静电动辄几千伏，一次触碰就可能击穿片内结构；但钳位也不是无限保险——它只为瞬态小能量设计，若把 12V 长期误接到 5V 引脚上，持续大电流会先烧掉钳位二极管，再烧掉芯片。

第三个角色是续流：给电感电流留一条退路。电感电流不能突变。继电器线圈由晶体管驱动，导通时电流 100mA；晶体管约在 1$\mu s$ 内关断，线圈仍要维持电流，两端感应电压可按 $V = L \cdot \Delta I / \Delta t$ 估算：10mH 电感在 1$\mu s$ 内切断 100mA，理想估算达 $10mH \times 100mA/1\mu s = 1000V$，实际虽被寄生电容和器件击穿限制，仍常达几十到上百伏，远超小信号晶体管的耐压。接法是在线圈两端反并一只二极管：正常供电时它反偏不工作，关断瞬间线圈电流改经二极管环流、慢慢衰减，反峰被钳在约 0.7V。没有续流二极管的继电器驱动，常见现象是晶体管工作一段时间后“莫名其妙”损坏。

\begin{longtable}{L{0.18\linewidth}L{0.4\linewidth}L{0.32\linewidth}}
\toprule
角色 & 典型接法 & 缺少它的失效后果 \\
\midrule
整流与防反接 & 单管串联或四管桥，串在电源入口 & 反向或交变电压直达直流母线，电容与稳压器受损 \\\tablerowrule
钳位保护 & 两只二极管把引脚反并到电源轨和地 & 静电与瞬态直接打入片内结构；长期过压同样会烧毁钳位 \\\tablerowrule
续流 & 与继电器或电感线圈反并联 & 关断反峰击穿驱动晶体管，表现为多次动作后失效 \\
\bottomrule
\end{longtable}

这三个角色都不参与逻辑判断，却决定了逻辑电路能否在真实电源和真实负载下长期存活。

二极管逻辑可以表达某些简单关系。若有两路输入想共同点亮一个指示节点，任意一路为高时都能通过二极管把输出节点拉高，输出就呈现 OR 行为：

\begin{lstlisting}
// 注释（推进）：先判断条件，再只执行命中的分支；分支顺序同时表达优先级。
// 注释（边界）：所有输入稳定后才读取结果；未覆盖的条件必须有明确默认行为。
if A can drive through a diode:
  Y may be pulled high
if B can drive through a diode:
  Y may be pulled high
\end{lstlisting}

该例能够说明方向性逻辑，但也暴露出局限。第一，二极管会带来压降，输出高电平可能比输入高电平低一截；级联多了，电平余量会越来越紧。第二，二极管只能提供方向性，不能主动把一个偏弱的高电平重新恢复成强高电平。第三，它不方便实现取反，而取反是构造通用逻辑必须具备的能力。

\begin{circuit}[x=1cm,y=1cm]
  % A 路二极管
  \node[hw label] at (0,1.4) {A};
  \draw[line width=0.55pt] (0.3,1.4) -- (0.7,1.4);
  \draw[line width=0.55pt, fill=white] (0.7,1.58) -- (1.06,1.4) -- (0.7,1.22) -- cycle;
  \draw[line width=0.55pt] (1.06,1.58) -- (1.06,1.22);
  \draw[line width=0.55pt] (1.06,1.4) -- (2.0,1.4) -- (2.0,0.75);
  % B 路二极管
  \node[hw label] at (0,-1.4) {B};
  \draw[line width=0.55pt] (0.3,-1.4) -- (0.7,-1.4);
  \draw[line width=0.55pt, fill=white] (0.7,-1.22) -- (1.06,-1.4) -- (0.7,-1.58) -- cycle;
  \draw[line width=0.55pt] (1.06,-1.22) -- (1.06,-1.58);
  \draw[line width=0.55pt] (1.06,-1.4) -- (2.0,-1.4) -- (2.0,-0.75);
  % 汇合节点与输出
  \draw[line width=0.55pt] (2.0,0.75) -- (2.0,-0.15);
  \draw[fill=BookInk] (2.0,0) circle (0.05);
  \draw[line width=0.55pt] (2.0,0) -- (3.4,0);
  \node[hw label] at (3.4,0.35) {Y};
  % 下拉电阻
  \draw[line width=0.55pt] (2.0,-0.15) -- (2.0,-0.32);
  \draw[line width=0.55pt] (1.82,-0.32) rectangle (2.18,-1.0);
  \draw[line width=0.55pt] (2.0,-1.0) -- (2.0,-1.15);
  \pic at (2.0,-1.15) {hw gnd};
  \node[hw label] at (2.95,-0.68) {下拉电阻};
  % 压降标注
  \node[hw label, text=BookAmber, text width=4.6cm] at (5.6,0.75) {A=5V 导通时\\$Y \approx 5.0-0.7 = 4.3$V};
  \node[hw label, text width=4.6cm] at (5.6,-0.75) {两路都低：\\Y 靠下拉电阻给出默认 0};
\end{circuit}
\diagramnote{二极管 OR 的真实样子：任意一路为高，该路二极管导通，把 Y 拉高——但 Y 比输入低一个正向压降（硅管约 0.7V）；两路都低时没有驱动源，靠下拉电阻给出默认 0。方向性由二极管保证，电平却不会被恢复：这就是它不能无限级联的原因。}

以安全启动控制板的故障汇总为例，\code{temp_alarm} 和 \code{current_alarm} 都可能点亮 \code{fault_lamp}。二极管 OR 可以让任一报警输入把故障节点拉高，但如果这个故障节点还要继续驱动后面的禁止逻辑、记录逻辑和指示灯，问题就会扩大。每经过一级二极管，电平余量都可能减少一部分；每增加一个负载，边沿都会变慢。最终出现的不是“公式错了”，而是高电平不够高、上升不够快、后级读不稳。

\begin{longtable}{L{0.24\linewidth}L{0.32\linewidth}L{0.34\linewidth}}
\toprule
二极管 OR 用法 & 可以解决的问题 & 不能解决的问题 \\
\midrule
汇总多个报警源 & 任一报警能提供一条拉高路径 & 不能主动恢复被削弱的高电平 \\\tablerowrule
隔离多个输入源 & 一个输入不容易反向灌入另一个输入 & 正向压降会消耗电平余量 \\\tablerowrule
形成简单有线关系 & 结构直观，适合早期理解 & 不适合无限级联成通用逻辑网络 \\
\bottomrule
\end{longtable}

取反为什么重要？因为许多条件天然需要反向解释。传感器未触发时允许运行，触发时禁止运行；按钮未按下时保持，按下时改变。若电路只能表达“有一路导通”，却不能把 1 变成 0、把 0 变成 1，那么可组合的逻辑很快就受限。要继续向前，需要一个信号能控制另一条电流路径是否导通。

晶体管的细节很深，但作为第一层硬件模型，可以先把它当成受控开关。控制端不是直接给输出供电，而是决定输出节点到某条电源路径之间是否接通。受控开关和普通机械开关最大的差别，是控制信号本身也是电信号。按钮需要人按，晶体管可以由前一级逻辑输出控制。于是硬件有了“信号控制信号”的递归能力：一个门的输出可以控制下一批开关，下一批开关再形成新的输出。

在 CMOS 数字电路中，可以先用极简模型区分两类受控开关。NMOS 更适合把节点拉向地，输入控制为高时容易导通；PMOS 更适合把节点拉向电源，输入控制为低时容易导通。真实器件还有阈值电压、体效应、漏电、沟道电阻和寄生电容等细节，但“数字硬件基础”部分暂时只抓住一件事：上拉路径和下拉路径由互补器件协作，才能把输出恢复为强 0 或强 1。

\begin{longtable}{L{0.18\linewidth}L{0.32\linewidth}L{0.4\linewidth}}
\toprule
器件视角 & 在门电路中的典型职责 & 需要避免的误解 \\
\midrule
NMOS 下拉 & 给输出提供到地的受控路径 & 不是“产生 0”这个符号，而是放电路径 \\\tablerowrule
PMOS 上拉 & 给输出提供到电源的受控路径 & 不是“产生 1”这个符号，而是充电路径 \\\tablerowrule
互补配对 & 一个负责拉低，另一个负责拉高 & 不能让两条强路径长期同时导通 \\\tablerowrule
输入过渡 & 两类器件可能都部分导通 & 过渡区只应短暂经过，不应长期停留 \\
\bottomrule
\end{longtable}

反相器可以作为第一种完整逻辑门来理解。输入低时，上拉路径导通、下拉路径断开，输出被拉高；输入高时，上拉路径断开、下拉路径导通，输出被拉低。注意这里不是“输入低所以输出自动高”，而是“输入低使某个高电平路径导通”。把每一句逻辑描述翻译成路径描述，能避免很多硬件误解。

\begin{longtable}{L{0.18\linewidth}L{0.29\linewidth}L{0.29\linewidth}L{0.14\linewidth}}
\toprule
输入 & 上拉路径 & 下拉路径 & 输出 \\
\midrule
低电平 & 导通 & 断开 & 被拉高 \\\tablerowrule
高电平 & 断开 & 导通 & 被拉低 \\\tablerowrule
过渡中 & 可能部分导通 & 可能部分导通 & 暂时不应采样 \\
\bottomrule
\end{longtable}

\begin{circuit}[x=1cm,y=1cm]
  \node[hw label] at (0,2.75) {VCC};
  \draw[line width=0.55pt] (0,2.6) -- (0,2.0);
  \pic at (0,1.58) {hw pmos};
  \draw[line width=0.55pt] (0,1.16) -- (0,0.84);
  \draw[fill=BookInk] (0,1.0) circle (0.05);
  \draw[line width=0.55pt] (0,1.0) -- (1.5,1.0);
  \node[hw label] at (1.5,1.35) {Y};
  \pic at (0,0.42) {hw nmos};
  \draw[line width=0.55pt] (0,0) -- (0,-0.15);
  \pic at (0,-0.15) {hw gnd};
  \node[hw label] at (-2.2,1.0) {A};
  \draw[line width=0.55pt] (-1.95,1.0) -- (-1.1,1.0);
  \draw[fill=BookInk] (-1.1,1.0) circle (0.05);
  \draw[line width=0.55pt] (-1.1,1.0) -- (-1.1,1.58) -- (-0.65,1.58);
  \draw[line width=0.55pt] (-1.1,1.0) -- (-1.1,0.42) -- (-0.65,0.42);
  \node[hw label, text=BookTeal, text width=3.4cm] at (3.6,1.75) {A=0：上通下断\\Y 被充电到高};
  \node[hw label, text=BookAmber, text width=3.4cm] at (3.6,0.15) {A=1：上断下通\\Y 被放电到低};
  \node[hw label] at (-0.75,2.0) {PMOS};
  \node[hw label] at (-0.75,0.0) {NMOS};
\end{circuit}
\diagramnote{反相器的开关模型。输入 A 同时接到两个控制端：PMOS 控制端带小圆圈，表示“输入低才导通”；NMOS 没有圆圈，表示“输入高才导通”。这就是 NOT 的全部——但它同时完成了二极管做不到的事：不管输入边沿多差，输出都被重新驱动成强 0 或强 1，电平被恢复了。}

把这张路径表压缩成真值表，就是 NOT 的全部定义：

\begin{longtable}{L{0.14\linewidth}L{0.14\linewidth}L{0.6\linewidth}}
\toprule
A & Y & 路径解释 \\
\midrule
0 & 1 & 上拉导通，输出被拉到高 \\\tablerowrule
1 & 0 & 下拉导通，输出被拉到低 \\
\bottomrule
\end{longtable}

第三行很重要。真实输入从低到高不是瞬间跳变，在过渡区间里，上拉和下拉可能都不处于理想开/关状态。短时间内出现电源到地的电流并不奇怪，这就是动态功耗和短路功耗的一部分。只要过渡足够快、采样避开中间区域，数字抽象仍然成立。若输入边沿太慢，电路可能在中间区域停留过久，功耗和误判风险都会上升。

把开关串联，可以表达“两个条件都满足才有通路”；把开关并联，可以表达“任意一个条件满足就有通路”。以下拉网络为例，两个受控开关串联，只有 A 和 B 都让开关导通时，输出节点才会被拉向低电平；两个开关并联，只要 A 或 B 有一个导通，输出节点就会被拉低。上拉网络通常要和下拉网络互补，保证输出在稳定输入下要么被拉高，要么被拉低，而不是两边都断或两边都强通。这种“互补路径”思想，后面会自然走到 CMOS 逻辑。

互补路径可以用两条规则记住。第一，输出应该被明确驱动：下拉网络不导通时，上拉网络应当负责把输出拉高；下拉网络导通时，上拉网络应当关闭，避免电源到地的强冲突。第二，稳定输入下不应依赖中间电压：如果上拉和下拉都处于部分导通状态，输出可能落在不可靠区间，功耗也会上升。CMOS 逻辑之所以适合作为数字门的基础，关键就在于它把“拉高”和“拉低”组织成互补关系，并在稳定状态下尽量避免静态直通电流。

\begin{longtable}{L{0.22\linewidth}L{0.34\linewidth}L{0.34\linewidth}}
\toprule
路径状态 & 输出后果 & 设计判断 \\
\midrule
上拉开、下拉关 & 输出被拉向高电平 & 合法稳定状态 \\\tablerowrule
上拉关、下拉开 & 输出被拉向低电平 & 合法稳定状态 \\\tablerowrule
上拉关、下拉关 & 输出悬空，可能保留旧电压 & 需要补默认路径或改逻辑 \\\tablerowrule
上拉开、下拉开 & 电源到地形成强通路 & 功耗高，电平不可靠 \\\tablerowrule
两边部分导通 & 输出可能处于过渡区 & 只允许短时间经过，不应长期停留 \\
\bottomrule
\end{longtable}

把这条规则应用到 NAND 和 NOR，可以看到 CMOS 门并不是把 AND、OR 符号直接画在硅片上，而是把上拉网络和下拉网络组织成互补关系。二输入 NAND 的下拉网络由两个 NMOS 串联：只有 A、B 都为 1 时，输出才被拉向地；它的上拉网络由两个 PMOS 并联：只要 A 或 B 有一个为 0，就有一条上拉路径把输出拉高。二输入 NOR 则相反：下拉网络由两个 NMOS 并联，只要 A 或 B 为 1 就能拉低；上拉网络由两个 PMOS 串联，只有 A、B 都为 0 时才拉高。

\begin{longtable}{L{0.16\linewidth}L{0.34\linewidth}L{0.4\linewidth}}
\toprule
CMOS 门 & 下拉网络 & 上拉网络 \\
\midrule
NAND & NMOS 串联，A=1 且 B=1 时完整导通 & PMOS 并联，A=0 或 B=0 时至少一路导通 \\\tablerowrule
NOR & NMOS 并联，A=1 或 B=1 时至少一路导通 & PMOS 串联，A=0 且 B=0 时完整导通 \\
\bottomrule
\end{longtable}

\begin{circuit}[x=1cm,y=1cm]
  \node[hw label] at (0,3.5) {VCC};
  \draw[line width=0.55pt] (-0.55,3.3) -- (-0.55,3.05);
  \draw[line width=0.55pt] (0.55,3.3) -- (0.55,3.05);
  \draw[line width=0.55pt] (-0.55,3.3) -- (0.55,3.3);
  \pic at (-0.55,2.63) {hw pmos};
  \pic at (0.55,2.63) {hw pmos};
  \draw[line width=0.55pt] (-0.55,2.21) -- (0,2.21) -- (0.55,2.21);
  \draw[line width=0.55pt] (0,2.21) -- (0,1.9);
  \draw[fill=BookInk] (0,2.05) circle (0.05);
  \draw[line width=0.55pt] (0,2.05) -- (2.2,2.05);
  \node[hw label] at (2.2,2.4) {Y};
  \pic at (0,1.48) {hw nmos};
  \pic at (0,0.62) {hw nmos};
  \draw[line width=0.55pt] (0,1.06) -- (0,1.04);
  \draw[line width=0.55pt] (0,0.2) -- (0,0.0);
  \pic at (0,0.0) {hw gnd};
  % A 输入：接 PMOS-A 与 NMOS-A
  \node[hw label] at (-2.9,1.48) {A};
  \draw[line width=0.55pt] (-2.65,1.48) -- (-0.65,1.48);
  \draw[fill=BookInk] (-1.9,1.48) circle (0.05);
  \draw[line width=0.55pt] (-1.9,1.48) -- (-1.9,2.63) -- (-1.2,2.63);
  % B 输入：接 PMOS-B 与 NMOS-B
  \node[hw label] at (-2.9,0.62) {B};
  \draw[line width=0.55pt] (-2.65,0.62) -- (-0.65,0.62);
  \draw[fill=BookInk] (-1.5,0.62) circle (0.05);
  \draw[line width=0.55pt] (-1.5,0.62) -- (-1.5,2.63) -- (-0.1,2.63);
  \node[hw label, text width=3.5cm] at (4.2,2.9) {上拉网络\\两个 PMOS 并联\\A=0 或 B=0 即拉高};
  \node[hw label, text width=3.5cm] at (4.2,0.55) {下拉网络\\两个 NMOS 串联\\A=1 且 B=1 才拉低};
\end{circuit}
\diagramnote{二输入 NAND 的晶体管级结构。输出为 0 的唯一情形是 A=B=1：两个串联 NMOS 同时导通，下拉路径完整；其余三行至少一只 PMOS 导通，上拉把输出充到高。真值表的每一行都能在这张图上找到对应的路径状态。（输入 B 的连线跨过 A 的线但不连接：没有连接点的交叉不算接通。）}

NOR 用同一套符号画出来，两个网络的串并联关系正好与 NAND 对调：

\begin{circuit}[x=1cm,y=1cm]
  \node[hw label] at (0,3.5) {VCC};
  \draw[line width=0.55pt] (-0.3,3.3) -- (0.3,3.3);
  \draw[line width=0.55pt] (0,3.3) -- (0,3.05);
  \pic at (0,2.63) {hw pmos};
  \draw[line width=0.55pt] (0,2.21) -- (0,2.05);
  \pic at (0,1.63) {hw pmos};
  \draw[line width=0.55pt] (0,1.21) -- (0,1.05);
  \draw[fill=BookInk] (0,1.05) circle (0.05);
  \draw[line width=0.55pt] (0,1.05) -- (2.2,1.05);
  \node[hw label] at (2.2,1.4) {Y};
  \draw[line width=0.55pt] (0,1.05) -- (0,0.77);
  \draw[line width=0.55pt] (-0.55,0.77) -- (0.55,0.77);
  \pic at (-0.55,0.35) {hw nmos};
  \pic at (0.55,0.35) {hw nmos};
  \draw[line width=0.55pt] (-0.55,-0.07) -- (-0.55,-0.45);
  \draw[line width=0.55pt] (0.55,-0.07) -- (0.55,-0.45);
  \pic at (-0.55,-0.45) {hw gnd};
  \pic at (0.55,-0.45) {hw gnd};
  % A 输入：接左 NMOS 与串联下管 PMOS
  \node[hw label] at (-2.9,0.35) {A};
  \draw[line width=0.55pt] (-2.65,0.35) -- (-1.25,0.35);
  \draw[fill=BookInk] (-1.9,0.35) circle (0.05);
  \draw[line width=0.55pt] (-1.9,0.35) -- (-1.9,1.63) -- (-0.7,1.63);
  % B 输入：接串联上管 PMOS 与右 NMOS（跨过输出线但不连接）
  \node[hw label] at (-2.9,2.63) {B};
  \draw[line width=0.55pt] (-2.65,2.63) -- (-0.7,2.63);
  \draw[fill=BookInk] (-0.45,2.63) circle (0.05);
  \draw[line width=0.55pt] (-0.45,2.63) -- (-0.45,3.0) -- (1.5,3.0) -- (1.5,0.35) -- (0.5,0.35);
  \node[hw label, text width=3.5cm] at (4.2,2.6) {上拉网络\\两个 PMOS 串联\\A=0 且 B=0 才拉高};
  \node[hw label, text width=3.5cm] at (4.2,0.35) {下拉网络\\两个 NMOS 并联\\A=1 或 B=1 即拉低};
\end{circuit}
\diagramnote{二输入 NOR 的晶体管级结构，与上一张 NAND 图成对：上拉网络换成两个 PMOS 串联，只有 A=B=0 时上拉路径才完整、输出被拉到高；下拉网络换成两个 NMOS 并联，A 或 B 任一为 1 就有一条下拉路径把输出拉低。输出为 1 的唯一稳定条件对应“上拉串联路径完整”，与互补网络表逐行对应。（输入 B 的连线跨过输出线但不连接：没有连接点的交叉不算接通。）}

这张表把逻辑表和器件路径接在一起。NAND 输出为 0 的唯一稳定条件，对应“下拉串联路径完整”；NAND 输出为 1 的三个稳定条件，对应“至少一条上拉路径存在”。NOR 输出为 1 的唯一稳定条件，对应“上拉串联路径完整”；NOR 输出为 0 的三个稳定条件，对应“至少一条下拉路径存在”。如果只背真值表，读者很难解释为什么 NAND/NOR 常成为门级实现的基础；如果只看器件路径，又容易失去逻辑函数的目的。两者必须同时保留。

\begin{classicnote}
CSAPP 一类系统教材常把“布尔函数”和“硬件实现”分层书写：先说明抽象函数，再说明门级和时序代价。这里的 CMOS NAND/NOR 正好提供一个小型范例。稳定逻辑值来自互补路径，传播延迟来自器件和连线，功耗来自节点切换和漏电；同一个门符号背后同时存在功能、时间和能量三个维度。
\end{classicnote}

从功耗角度看，数字门并非“只有 0 和 1，没有模拟代价”。稳定状态下，如果只有上拉或只有下拉导通，静态电流可以很小；切换时，输出节点的寄生电容需要充电或放电，形成动态功耗；输入经过过渡区时，上拉和下拉可能短时间同时导通，形成短路功耗。

这也是 CMOS 成为主流数字逻辑基础的重要原因。它不是没有功耗，而是在稳定 0 和稳定 1 时尽量避免持续直通电流，把主要能量消耗集中到节点切换和短暂过渡中。随着芯片集成度提高，设计者更关心哪些节点频繁翻转、哪些长线负载过大、哪些门需要更强驱动，以及哪些路径应由寄存器边界切分。

\begin{longtable}{L{0.2\linewidth}L{0.34\linewidth}L{0.36\linewidth}}
\toprule
功耗来源 & 物理原因 & 设计含义 \\
\midrule
静态功耗 & 漏电或稳定状态下仍存在电流路径 & 器件越小越不能忽略漏电和待机功耗 \\\tablerowrule
动态功耗 & 输出节点和连线电容反复充放电 & 高翻转率、长线和大扇出都会增加能耗 \\\tablerowrule
短路功耗 & 输入过渡时上拉和下拉短暂同时导通 & 慢边沿不仅影响时序，也会增加功耗 \\\tablerowrule
驱动损耗 & 输出级推动外部负载或长线 & 指示灯、总线和连接器常需要专用驱动级 \\
\bottomrule
\end{longtable}

\topic{BJT 与 MOSFET 开关对照}
本章的受控开关主要以 MOS 管出现，但值得把它和另一类器件——双极型晶体管（BJT）——放在一起对照，因为两者的差别正好解释了现代数字电路的工艺选择。最核心的差别在控制端：BJT 是电流驱动器件，集电极电流要靠持续的基极电流维持，工程估算常取 $I_B \ge I_C/\beta$，$\beta$ 约几十到几百；MOSFET 是电压驱动器件，栅极与沟道之间隔着绝缘氧化层，稳定状态下几乎没有栅极电流，能量只在翻转瞬间用于给栅电容充放电。

作开关使用时，两者的导通损耗来源也不同。BJT 在截止与饱和之间切换，导通时呈现近似恒定的饱和压降 $V_{CE(sat)}$（约 0.2V），损耗为压降乘以电流；MOSFET 在夹断与导通之间切换，导通时表现为沟道电阻 $R_{DS(on)}$，损耗按 $I^2R$ 计算。算一笔账：驱动一个 100mA 的继电器，BJT 取 $\beta=100$，基极要持续供出约 1mA，导通损耗约 $0.2V \times 100mA = 20mW$；换用 $R_{DS(on)}=50m\Omega$ 的 MOSFET，损耗只有 $0.1^2 \times 0.05 = 0.5mW$，且栅极在稳定期间几乎不取电流。

\begin{longtable}{L{0.2\linewidth}L{0.36\linewidth}L{0.34\linewidth}}
\toprule
对比项 & BJT & MOSFET \\
\midrule
驱动方式 & 电流驱动，基极需持续电流 & 电压驱动，栅极只取充放电电流 \\\tablerowrule
开关状态 & 截止与饱和 & 夹断与导通 \\\tablerowrule
导通损耗来源 & $V_{CE(sat)} \times I$，近似恒定压降 & $I^2 \times R_{DS(on)}$，随电流平方上升 \\\tablerowrule
100mA 驱动例 & 基极约 1mA，损耗约 20mW & 栅极静态近零，损耗约 0.5mW（50m$\Omega$） \\\tablerowrule
静态功耗 & 基极电流贯穿整个导通期 & 稳定状态下几乎为零 \\\tablerowrule
仍活跃的场合 & 大电流驱动、线性放大、分立高压 & 数字集成与功率开关的主流 \\
\bottomrule
\end{longtable}

由此可以理解为什么现代数字电路主流是 MOS。数字芯片里有数十亿个开关，若每个开关的控制端都要持续电流，静态功耗将完全无法接受；MOS 栅极不取静态电流，NMOS 与 PMOS 两种极性又容易在同一工艺中做成互补对，器件尺寸还能持续缩小——静态功耗和集成密度因此都站在 MOS 一边。BJT 并未消失：在大电流驱动、线性放大和某些高压分立场合，它的线性区特性和电流处理能力仍然合适，只是不再是构成通用数字逻辑的器件。

开漏或开集结构是理解共享线路的重要例子。一个输出器件只负责把线拉低，释放时并不主动拉高；高电平由公共上拉电阻提供。多个器件可以共享同一根线，只要任意一个器件拉低，整根线就是低电平；只有所有器件都释放，线上才被上拉为高电平。这种结构常用于中断请求、故障汇总或简单总线。

\begin{longtable}{L{0.22\linewidth}L{0.34\linewidth}L{0.34\linewidth}}
\toprule
共享线状态 & 物理路径 & 逻辑后果 \\
\midrule
所有输出释放 & 上拉电阻把线拉高 & 线为 1，但上升速度取决于上拉和负载 \\\tablerowrule
任一输出拉低 & 该输出提供到地路径 & 线为 0，多个输出不会互相强冲突 \\\tablerowrule
上拉过弱 & 高电平上升慢 & 采样前可能尚未达到高阈值 \\\tablerowrule
上拉过强 & 拉低时电流大 & 功耗增加，输出器件负担加重 \\
\bottomrule
\end{longtable}

这说明“多个来源接到一根线”并非永远错误，关键在于是否有协议和电气结构保证同一时刻不会出现强驱动冲突。开漏共享线允许多个设备共同拉低，是因为它们只竞争“是否提供到地路径”，不竞争“谁主动拉高”。普通推挽输出若直接并在一起，一个输出拉高、另一个输出拉低，就会形成强冲突。

三态输出是另一类共享线路方案。普通数字输出只有主动拉高和主动拉低两种驱动状态；三态输出还可以进入高阻态。高阻态不是逻辑 0，也不是逻辑 1，而是输出驱动器基本断开，让这一路暂时不参与驱动总线。若多块设备共享数据总线，控制逻辑必须保证同一时刻最多只有一个设备打开输出驱动，其余设备都处于高阻态。否则，一个设备驱动 1、另一个设备驱动 0，仍会发生总线争用。

\begin{longtable}{L{0.18\linewidth}L{0.28\linewidth}L{0.44\linewidth}}
\toprule
输出方式 & 驱动状态 & 适合表达的问题 \\
\midrule
推挽输出 & 主动拉高或主动拉低 & 单一来源驱动一条内部信号 \\\tablerowrule
开漏输出 & 只主动拉低，释放靠上拉 & 多个来源共同提出低有效请求 \\\tablerowrule
三态输出 & 拉高、拉低或高阻 & 多个来源按仲裁结果轮流占用总线 \\
\bottomrule
\end{longtable}

开漏和三态都要求协议。开漏协议通常是“任一设备可以拉低，所有设备释放后才为高”；三态协议通常是“仲裁者只允许一个设备驱动，其他设备必须释放”。这两类结构后来会在中断线、片选、数据总线和 I/O 控制中反复出现。现在先记住硬件判断标准：多来源共享一根线时，必须说明每个来源什么时候能驱动、什么时候必须释放，以及释放后由谁给出默认电平。

数字门不是画在纸上的符号，而是一组受输入控制的上拉、下拉、释放和恢复路径。二极管提供方向性，但不恢复电平；晶体管提供受控路径；互补开关网络提供明确的高低驱动；缓冲和三态驱动器提供负载隔离和共享边界。

器件族之间还存在电平兼容问题。即使两个器件都叫“数字逻辑”，也不表示它们可以任意直连。一个 5V 供电器件的输出、一个 3.3V CMOS 输入、一个 TTL 兼容输入和一个开漏传感器，可能有不同的高低阈值、输入电流和输出驱动能力。设计者必须把数据手册中的最坏输出和最坏输入放在同一张表里比较，而不是只看“都是 0/1”。

\begin{longtable}{L{0.22\linewidth}L{0.34\linewidth}L{0.34\linewidth}}
\toprule
连接检查 & 要问的问题 & 典型风险 \\
\midrule
电源电压 & 前级输出高电平是否超过后级输入上限 & 过压损坏或长期可靠性下降 \\\tablerowrule
输入阈值 & 前级最差高/低输出是否满足后级阈值 & 电平看似变化，后级不承认 \\\tablerowrule
输入电流 & 前级能否提供或吸收后级所需电流 & 电平被拖偏，余量不足 \\\tablerowrule
输出驱动 & 前级能否驱动全部负载和走线电容 & 边沿变慢，时序失败 \\\tablerowrule
上拉配置 & 开漏/开集输出是否有合适上拉 & 释放后没有可靠高电平 \\
\bottomrule
\end{longtable}

\topic{5V 与 3.3V 系统互连的电平转换}
上表回答的是“怎么查”，本专题回答“查出不匹配之后怎么办”。先把两个方向各自卡在哪里算清楚。5V 输出进 3.3V 输入：5V 超过多数 3.3V 器件的最大输入额定值，长期直连可能损坏输入结构，除非该输入明确标注 5V 容忍。反方向的问题更隐蔽：3.3V 系统的最差高电平可能只有 2.4V 左右，而 74HC 类 CMOS 输入在 5V 供电时要求 $V_{IH} \ge 0.7V_{CC} = 3.5V$，电平看似翻转了，后级却不承认。同一 5V 供电下，74HCT 系列把输入阈值做成 TTL 兼容：$V_{IH}=2.0V$、$V_{IL}=0.8V$，2.4V 的高电平仍有 0.4V 余量。74HC 与 74HCT 的分工由此而来：同一电源域内部级联用 HC，需要承接 TTL 电平或 3.3V 输出的场合用 HCT。

真正需要转换时，常见做法有三种，代价各不相同。第一种是电阻分压：上臂 1.7k$\Omega$、下臂 3.3k$\Omega$，可把 5V 分成 $5V \times 3.3/(1.7+3.3) = 3.3V$。它只适合从高到低、单向、低速的信号：分压网络的等效内阻（本例约 1.1k$\Omega$）会拖慢边沿，$5V/5k\Omega=1mA$ 的静态电流常流不断，而且电阻只会衰减不会放大，3.3V 到 5V 方向完全用不了。第二种是专用电平转换芯片：双向、多通道，阈值和速度在数据手册中有明确规格；代价是多一片器件、两组电源脚和相应的布板面积。第三种是开漏加目标电源上拉：输出级只做开漏，上拉电阻接到对方电源——器件只负责拉低，高电平天然由目标电源域提供，电平自动匹配；I2C 总线正是靠这个结构允许不同电压域挂在同一条线上。代价是速度受上拉 RC 限制，且只有开漏或开集结构的信号才能这样接。

\begin{longtable}{L{0.26\linewidth}L{0.32\linewidth}L{0.32\linewidth}}
\toprule
做法 & 适用方向与速度 & 代价与限制 \\
\midrule
电阻分压（1.7k$\Omega$+3.3k$\Omega$） & 仅 5V 到 3.3V、单向、低速 & 约 1mA 静态电流，约 1.1k$\Omega$ 等效内阻拖慢边沿，不能反向 \\\tablerowrule
电平转换芯片 & 双向、多通道，速度有手册规格 & 增加器件、两组电源和布板面积 \\\tablerowrule
开漏加目标域上拉 & 双向共享线，适合总线与中断请求 & 速度受上拉 RC 限制，仅开漏结构可用 \\\tablerowrule
选对输入阈值（74HC 与 74HCT） & 3.3V 到 5V 单方向最省事 & 只是选型手段，不解决 5V 到 3.3V 过压 \\
\bottomrule
\end{longtable}

三种做法都不改变逻辑内容，只重新匹配两侧的电平约定。选型顺序是先定方向、速度和通道数，再决定用分压、专用芯片还是开漏结构。

数据手册里的参数要按“最坏组合”阅读。若前级在轻载、常温下能输出漂亮的 3.3V，不代表它在最大负载、高温、低电源电压下仍有同样余量；若后级典型阈值较低，也不代表所有批次、所有温度下都按典型值判断。经典教材强调抽象层次，但并不鼓励忽略边界条件。抽象的意义是把边界写清楚，然后在边界内使用更高层概念。

读取标准逻辑芯片数据手册时，可以采用固定顺序。先看供电范围，确认器件是否处在允许电源电压内；再看输入阈值，确认前级输出能被后级可靠承认；再看输出电流能力，确认它能驱动实际后级和上拉/下拉网络；再看传播延迟，确认多级组合路径能在采样前稳定；最后看封装、引脚、温度范围和绝对最大额定值，确认板级使用没有越界。这个顺序比直接查“某个门是什么功能”更接近硬件设计的实际检查顺序。

\begin{longtable}{L{0.18\linewidth}L{0.36\linewidth}L{0.36\linewidth}}
\toprule
读取顺序 & 要确认的参数 & 与本章概念的关系 \\
\midrule
供电条件 & 工作电压范围、绝对最大额定值 & 电平约定必须先处在器件允许边界内 \\\tablerowrule
输入条件 & $V_\mathrm{IH}$、$V_\mathrm{IL}$、输入电流、滞回类型 & 后级如何把真实电压承认为 0/1 \\\tablerowrule
输出条件 & $V_\mathrm{OH}$、$V_\mathrm{OL}$、输出源/灌电流 & 前级能否给出足够强的高低电平 \\\tablerowrule
时间条件 & 传播延迟、上升/下降时间、使能/禁止延迟 & 组合路径多久后才可被使用 \\\tablerowrule
负载条件 & 扇出、输出电容、推荐上拉、总线负载 & 真实连接是否拖慢边沿或压缩余量 \\\tablerowrule
环境和封装 & 温度范围、封装引脚、ESD 和焊接约束 & 板级实现是否仍满足器件边界 \\
\bottomrule
\end{longtable}

例如同样是“用一个 NAND 门”，74LS、74HC、4000 CMOS 和高速 CMOS 的电气约定并不相同。某个系列可能适合 5V TTL 兼容输入，另一个系列可能要求 CMOS 阈值，某些高速器件对输入边沿和去耦要求更严格。正式设计不能只写“接一个 NAND”，还要写清器件族、供电电压、输入输出阈值、负载和时序要求。否则，逻辑表达式相同，实物板卡仍可能在温度、批次或负载变化时失效。

从这个角度看，缓冲器、电平转换器和驱动器并不是“逻辑上多余”的器件。缓冲器可以恢复电平和增强驱动；电平转换器让两个电源域的约定重新匹配；驱动器让小逻辑信号控制更大的负载。它们是否改变真值表不是唯一问题；它们是否让真实节点满足电平、负载和时序约定，才是硬件设计要回答的问题。

标准逻辑芯片把这些约束以产品形式暴露出来。以常见的 7400 系列为例，一个封装中可以放入若干个 NAND 门。教材里的 NAND 真值表只说明稳定逻辑行为，而芯片数据手册还会说明每个输入能承认的高低电平、输出能提供或吸收的电流、传播延迟范围、工作电压范围和扇出建议。换言之，标准门芯片把“布尔函数”和“电气约定”绑定在一起出售。

\begin{longtable}{L{0.2\linewidth}L{0.34\linewidth}L{0.36\linewidth}}
\toprule
从 7400 类门芯片读到的项目 & 对应本章概念 & 设计时的含义 \\
\midrule
输入高低阈值 & 可靠 bit & 前级输出必须落入后级承认区间 \\\tablerowrule
输出电流能力 & 驱动能力和扇出 & 一个门不能无限驱动后级 \\\tablerowrule
传播延迟 & 组合路径时序 & 多级门相连后要累计最坏延迟 \\\tablerowrule
工作电压范围 & 电平约定边界 & 不同电源域不能随意直连 \\\tablerowrule
封装引脚 & 物理边界 & 逻辑门最终仍要通过真实引脚连接 \\
\bottomrule
\end{longtable}

4000 系列 CMOS 逻辑则更适合说明互补路径和功耗。CMOS 门在稳定 0 或稳定 1 时，理想情况下没有从电源到地的持续强通路，静态功耗可以很低；切换时，输出电容充放电，形成动态功耗。现代处理器虽然远比 4000 系列复杂，但仍然继承了这个核心事实：频率越高、切换节点越多、负载越大，动态功耗越高；门越小、线越短、负载越低，延迟越容易下降。

比较 7400 TTL 和 4000 CMOS 时，不应把它们只看成“都能实现 NAND/NOR”。TTL 器件常用于说明输入电流、扇出和标准逻辑电平；CMOS 器件更适合说明极高输入阻抗、低静态功耗和输入悬空风险。一个 CMOS 输入若没有明确上拉或下拉，可能因为极小漏电和耦合电荷保持在不确定电平；这与“输入阻抗高”并不矛盾，反而正是高阻输入必须定义默认值的原因。

\begin{longtable}{L{0.2\linewidth}L{0.34\linewidth}L{0.36\linewidth}}
\toprule
芯片族视角 & 教材中应观察的事实 & 对本章的补充意义 \\
\midrule
7400 TTL & 输入电流、输出灌/拉电流、传播延迟和扇出限制 & 门芯片是带电气约定的布尔函数 \\\tablerowrule
4000 CMOS & 高输入阻抗、低静态功耗、宽电源范围和较慢边沿选型 & 默认路径和输入调理不能省略 \\\tablerowrule
高速 CMOS & 更低延迟、更强驱动、更严格边沿和布局要求 & 速度提高后，负载、布线和毛刺更敏感 \\
\bottomrule
\end{longtable}

\topic{体二极管与栅极保护}
高输入阻抗这条性质还有更深的器件根源，值得补两句。第一句关于体二极管。功率 MOSFET 的源极与衬底在制造时短接，衬底与漏极之间天然留下一个 PN 结，称为体二极管。对源极接地的 NMOS，它的方向是从源极指向漏极：正常工作时漏极电位高于源极，二极管反偏不导通；一旦漏极被拉到比源极低约 0.7V——例如走线电感上的负尖峰或电源接反——它就绕过栅极自行导通。栅极控制不了这只二极管，所以 MOSFET 开关挡不住反向电流；桥式驱动里它有时被有意用作续流路径，但长期依赖它导通会带来额外损耗和可靠性风险。

第二句关于栅氧击穿。MOS 栅极是一块被氧化层绝缘的电容极板，氧化层只有纳米量级厚，常见器件的栅源额定电压仅约 $\pm 20V$；而人体摩擦积累的静电轻松达到几千伏。栅极又是高阻节点，电荷没有泄放路径，一次静电触碰就可能把氧化层击穿，器件当场失效或留下隐患。这正是 CMOS 输入必须配保护结构、未用输入必须接确定电平、插拔与焊接要防静电的物理原因。

两个廉价元件能化解其中大部分风险。其一是在驱动输出与栅极之间串一只电阻：它与栅电容构成 RC，既限制瞬态进入栅极的电流，也阻尼快速边沿引起的振铃；取值按速度定，1k$\Omega$ 串阻配 5pF 小信号栅电容的时间常数约为 $1k\Omega \times 5pF = 5ns$，对低速输入毫无影响，驱动功率 MOS 时则常用几欧到几十欧换取边沿与开关损耗的平衡。其二是给栅极配一只到地或到电源的大阻值电阻，例如 100k$\Omega$：驱动断开、复位期间或连接器拔出时，栅极不再悬空，器件保持确定关断。这就是前面上拉/下拉原则在器件级的再次出现——高阻节点必须有默认归宿。

\begin{longtable}{L{0.2\linewidth}L{0.38\linewidth}L{0.32\linewidth}}
\toprule
风险 & 物理来源 & 常用对策 \\
\midrule
体二极管误导通 & 漏源电压反向超过约 0.7V & 抑制负压尖峰，必要时串肖特基二极管或按反峰设计 \\\tablerowrule
栅氧击穿 & 静电在无路可泄的高阻栅极上积累 & 输入保护结构、串联电阻、防静电操作 \\\tablerowrule
悬空误开通 & 驱动断开时栅极电荷随机 & 栅极上拉或下拉给出确定电平 \\\tablerowrule
边沿振铃 & 栅电容与走线电感形成谐振 & 串联栅电阻阻尼，按速度取值 \\
\bottomrule
\end{longtable}

这些措施都不改变真值表，却决定了开关器件在第一块板子上能否活着工作。

\begin{chipcase}
读一片标准 NAND 芯片，至少要同时读三层信息。第一层是逻辑：两个输入同时为 1 时输出为 0，其余情况输出为 1。第二层是电气：输入高低阈值、输出电流、工作电压和负载能力说明它能否与前后级可靠连接。第三层是时间：传播延迟说明输入变化后多久输出才可用。若只读第一层，就会把真实芯片误当成理想布尔符号；若只读第二、三层，又会失去门的功能目的。
\end{chipcase}

把逻辑式映射到芯片引脚时，还要完成一次\hwevidence{表达式到封装}的翻译。以 74HC00 或传统 7400 类四路二输入 NAND 芯片为例，一个封装中包含四个独立 NAND 门。设计者不能只写“使用一个 NAND”，还要指定哪两个信号接到某一路 NAND 的输入脚，该路输出脚驱动哪个后级，芯片的 \code{VCC} 和 \code{GND} 接到哪个电源域，电源脚附近是否放置去耦电容，未使用门的输入是否被接到确定电平。

\begin{longtable}{L{0.2\linewidth}L{0.36\linewidth}L{0.34\linewidth}}
\toprule
落地项目 & 需要写清的内容 & 常见错误 \\
\midrule
电源脚 & \code{VCC}、\code{GND}、工作电压和去耦电容位置 & 只接信号脚，忽略供电噪声和地参考 \\\tablerowrule
输入脚 & A、B 两个输入分别来自哪些内部语义信号 & 把低有效或未调理信号直接接入门 \\\tablerowrule
输出脚 & 输出驱动哪个门、缓冲器或后级输入 & 未检查扇出、负载和传播延迟 \\\tablerowrule
未用门 & 未使用输入接到确定高或低，输出可悬空不用 & CMOS 输入悬空导致随机翻转和额外功耗 \\\tablerowrule
封装方向 & 引脚编号、缺口方向、封装类型和焊接方向 & 原理图正确但实物引脚接反 \\
\bottomrule
\end{longtable}

例如要把 \code{safe_chain = cover_closed AND emergency_ok} 映射到 74HC00，可先用一路 NAND 得到反相信号，再用另一路 NAND 自反相恢复为 AND。文字规格可以写成：第 1 路 NAND 输入接 \code{cover_closed} 与 \code{emergency_ok}，输出命名为 \code{safe_chain_n}；第 2 路 NAND 两个输入都接 \code{safe_chain_n}，输出命名为 \code{safe_chain}。这种写法让中间反相信号、芯片资源占用和调试测点都可审查。

\begin{lstlisting}
safe_chain_n = cover_closed NAND emergency_ok
safe_chain   = safe_chain_n NAND safe_chain_n
\end{lstlisting}

\begin{pdfdiagram}
  \node[hw control, minimum width=2.0cm] (cover) at (0,0.75) {\code{cover\_closed}};
  \node[hw control, minimum width=2.0cm] (estop) at (0,-0.75) {\code{emergency\_ok}};
  \node[hw compute, minimum width=2.0cm] (nand1) at (3.1,0) {NAND\\第 1 路};
  \node[hw compute, minimum width=2.0cm] (nand2) at (6.2,0) {NAND\\第 2 路};
  \node[hw state, minimum width=2.25cm] (safe) at (9.4,0) {\code{safe\_chain}};
  \draw[hw bus] (cover.east) -- ++(0.55,0) |- (nand1.west);
  \draw[hw bus] (estop.east) -- ++(0.55,0) |- (nand1.west);
  \draw[hw bus] (nand1.east) -- (nand2.west);
  \node[hw label] at (5.15,1.25) {\code{safe\_chain\_n}};
  \draw[hw return] (nand1.east) -- ++(0.55,0.95) -| (nand2.north);
  \draw[hw return] (nand1.east) -- ++(0.55,-0.95) -| (nand2.south);
  \draw[hw bus] (nand2) -- (safe);
  \node[hw label, text width=9.6cm] at (4.7,-1.75) {第二路 NAND 的两个输入接同一个反相信号，等价于把第 1 路输出再反相。图中分成两条回入线，是为了避免箭头穿过文字。};
\end{pdfdiagram}
\diagramnote{用两路 NAND 实现 AND。图的命题是：功能等价必须能对应到具体芯片资源和中间测点。}

芯片引脚级说明还必须保留物理约束。每片芯片电源脚附近通常需要本地去耦电容，以降低门切换时的电源压降；输入不能超过允许电压范围；输出不能超出源电流或灌电流能力；同一输出驱动多个输入时要重新计算扇出；跨连接器或长线时还要考虑保护、终端和边沿质量。把布尔式映射到芯片，不是把公式抄到封装上，而是把功能、电平、时间、负载和装配边界同时写进规格。

若读者希望把安全启动控制板搭成一块低速演示电路，可以采用标准 CMOS 逻辑芯片完成第一版。这个版本不追求工业安全认证，只用于验证本章的电平、路径、组合逻辑和测量方法。输入侧使用按键或拨码开关表示 \code{start_button_raw} 等原始输入（分别对应 \code{start_request}、\code{cover_closed}、\code{emergency_ok}、\code{temp_alarm} 和 \code{current_alarm} 五个语义信号）；每个输入都有明确上拉或下拉，必要时经施密特输入整形；组合核心用 74HC00、74HC04、74HC08、74HC32 或同类芯片实现；输出侧用缓冲器或三极管/MOSFET 驱动 LED，而不是让普通门输出承担全部指示灯电流。

\begin{longtable}{L{0.22\linewidth}L{0.34\linewidth}L{0.34\linewidth}}
\toprule
模块 & 可选器件或结构 & 检查重点 \\
\midrule
输入默认值 & 10k$\Omega$ 量级上拉/下拉、电阻网络或拨码开关模块 & 松开、断线、未接时不悬空，默认方向符合安全策略 \\\tablerowrule
慢边沿整形 & 74HC14 或同类施密特反相器 & 慢按钮边沿和噪声不能在阈值附近反复翻转 \\\tablerowrule
基本门 & 74HC00、74HC08、74HC32、74HC04 或等价标准门 & 每个门的输入阈值、输出驱动和传播延迟满足连接要求 \\\tablerowrule
故障汇总 & OR 门、NAND-only 网络或二极管加缓冲结构 & 任一故障都点亮故障灯并禁止启动，双故障有明确行为 \\\tablerowrule
输出指示 & 74HC244/74HC125 缓冲、三极管或小 MOSFET 驱动 LED & 指示灯负载不拖垮内部逻辑电平 \\\tablerowrule
供电与去耦 & 5V 或 3.3V 稳定电源，每片芯片近旁 0.1$\mu F$ 去耦 & 切换时电源和地不出现影响阈值判断的扰动 \\
\bottomrule
\end{longtable}

搭建顺序应从静态边界开始，而不是直接运行全部逻辑。第一步只接电源、地和去耦，确认芯片供电正确。第二步只接输入默认网络，用万用表确认每个输入在松开、按下、断线时的电平。第三步逐个验证单门行为，例如 NAND、反相器和 OR 的四行真值表。第四步再连接 \code{fault_lamp} 和 \code{no_fault}，确认任一报警都会禁止启动。第五步连接 \code{safe_chain}、\code{request_ok} 和 \code{allow_start}，遍历正常允许、单条件缺失、单故障、双故障和断线默认。最后才连接 LED 驱动或执行机构模拟负载，并确认负载加入后内部逻辑电平仍满足后级阈值。

\begin{longtable}{L{0.14\linewidth}L{0.38\linewidth}L{0.38\linewidth}}
\toprule
阶段 & 操作 & 进入下一阶段的条件 \\
\midrule
1 & 只检查电源、地、去耦和芯片方向 & 供电电压正确，芯片不发热，电源脚附近无明显跌落 \\\tablerowrule
2 & 接入输入默认网络，不接组合核心 & 每个输入在所有机械状态下都有确定电平 \\\tablerowrule
3 & 单独验证每片门芯片的一路门 & 真值表正确，输出能驱动一个后级输入或测试点 \\\tablerowrule
4 & 接入故障汇总和无故障反相信号 & 报警输入任一为 1 时 \code{no_fault}=0 \\\tablerowrule
5 & 接入安全链和启动允许逻辑 & 只有所有允许条件同时成立时 \code{allow_start}=1 \\\tablerowrule
6 & 接入缓冲、LED 或模拟执行负载 & 负载不改变内部逻辑结果，边沿和电平仍可测量 \\
\bottomrule
\end{longtable}

这个低速演示电路也能暴露真实工程问题。若未用 CMOS 输入悬空，LED 可能随机闪烁；若按钮没有消抖，逻辑分析仪可能看到多个启动请求；若把 LED 直接接到弱输出上，门输出的高电平可能被拖低；若两片推挽输出误接到同一节点，一高一低时会产生冲突电流；若所有芯片共用长电源线而缺少去耦，多个输出同时翻转时可能出现误触发。它们都不是布尔代数错误，却都属于“数字硬件基础”部分必须掌握的硬件错误。

\section{三、案例：5V 外部状态怎样安全进入 3.3V 输入}

“高电平能够被识别”与“引脚能够长期承受该电压”是两个不同条件。设板外设备使用 5V 推挽输出，保证低电平不高于 0.4V、高电平不低于 4.4V；接收芯片工作在 3.3V，输入低阈值为 0.99V、输入高阈值为 2.31V，但绝对最大输入只有 $V_{\mathrm{DD}}+0.3\mathrm{V}$，约 3.6V。低电平同时满足功能与安全要求；高电平虽然远高于 2.31V，能够被判为 1，却超过了 3.6V 的绝对最大值，不能直接连接。

这组数字揭示了数据手册中的两类边界。$V_{\mathrm{IL}}$ 与 $V_{\mathrm{IH}}$ 说明接收器怎样解释电平，属于功能规格；绝对最大额定值说明器件不受损的极限，属于生存边界。绝对最大值不是推荐工作点，也不能因为串联电阻限制了电流，就默认可以持续依赖片内保护结构工作。

\begin{longtable}{L{0.20\linewidth}L{0.24\linewidth}L{0.25\linewidth}L{0.21\linewidth}}
\toprule
输入状态 & 功能判断 & 电气判断 & 结论 \\
\midrule
5V 端输出 0.4V & 低于 0.99V，可靠为 0 & 未超过负向额定范围 & 可作为合法低电平 \\\tablerowrule
5V 端输出 4.4V & 高于 2.31V，可靠为 1 & 高于约 3.6V 的绝对最大值 & 逻辑正确但电气不安全 \\\tablerowrule
接收端未上电 & 不能形成有效内部逻辑 & 高电平可能经钳位路径向 3.3V 电源反灌 & 必须阻断或限制反向供电 \\\tablerowrule
输入处于 0.99--2.31V & 不保证为 0 或 1 & 内部上下拉网络可能同时部分导通 & 不得作为稳定工作状态 \\
\bottomrule
\end{longtable}

\topic{串联电阻能限制异常电流，但不能自动完成电平转换}

若 5V 信号只是一次短暂异常，接收端钳位约在 3.6V，且数据手册允许钳位电流不超过 2mA，那么串联电阻的理论下限为
\[
  R_{\min}=\frac{5.0-3.6}{2\mathrm{mA}}=700\Omega.
\]
取 1k$\Omega$ 时，理想钳位电流约为 1.4mA。这个计算只回答“异常发生时电流是否受限”，没有证明片内钳位适合持续导通，也没有处理接收芯片掉电时的反灌。数据手册若没有明确允许这种工作方式，就应把它视为保护路径，而不是正常数据路径。

对于单向、低速状态信号，可以使用电阻分压。上臂 10k$\Omega$、下臂 20k$\Omega$ 时，5V 被分为约 3.33V；两只电阻的戴维南等效电阻约为 6.67k$\Omega$，还要与输入电容共同决定上升时间，并核对输入漏电在高温下造成的误差。对于高速信号、双向信号、接收端可能掉电的信号，使用明确支持 5V 输入且具备掉电隔离的缓冲器或电平转换器更稳妥。

\begin{pdfdiagram}[x=1cm,y=1cm]
  \node[hw state, minimum width=2.5cm] (source) at (1.4,4.0) {板外 5V 输出\\0.4V 或 4.4V};
  \node[hw control, minimum width=2.7cm] (protect) at (4.6,4.0) {连接器保护\\ESD/浪涌/限流};
  \node[hw compute, minimum width=2.8cm] (translate) at (8.0,4.0) {电平转换\\缓冲或电阻网络};
  \node[hw state, minimum width=2.6cm] (pin) at (11.3,4.0) {3.3V 输入脚\\安全电压范围};
  \node[hw control, minimum width=2.8cm] (threshold) at (11.3,1.2) {输入阈值\\低/未规定/高};
  \node[hw state, minimum width=2.8cm] (bit) at (8.0,1.2) {内部可靠 bit\\明确 0 或 1};
  \node[hw control, minimum width=2.7cm] (fault) at (4.6,1.2) {异常监测\\过压、断线、掉电};
  \node[hw state, minimum width=2.5cm] (safe) at (1.4,1.2) {安全结果\\隔离或默认状态};
  \draw[hw bus] (source) -- (protect);
  \draw[hw bus] (protect) -- (translate);
  \draw[hw bus] (translate) -- (pin);
  \draw[hw bus] (pin) -- (threshold);
  \draw[hw return] (threshold) -- (bit);
  \draw[hw return] (bit) -- (fault);
  \draw[hw return] (fault) -- (safe);
\end{pdfdiagram}
\diagramnote{5V 外部状态进入 3.3V 逻辑的完整边界。上排先限制异常能量并完成电平适配，下排才由输入阈值形成内部 bit；过压、断线或接收端掉电时，故障路径必须把输入隔离或带到规定默认状态，不能让片内钳位长期承担电平转换。}

\topic{掉电顺序决定是否出现反向供电}

接收芯片掉电时，$V_{\mathrm{DD}}$ 变为 0V。若外部 5V 输出仍保持高电平，输入钳位可能把电流送入已经关闭的 3.3V 电源网，使芯片局部处于半上电状态。此时内部复位、I/O 状态和功耗都不再受正常规格保证；电流还可能经其他引脚继续流向板上器件。具备 \term{掉电保护}{power-off protection} 或 \term{部分掉电模式}{partial power-down mode} 的缓冲器，会在自身或接收端未供电时保持高阻，切断这条路径。

调试时应同时测三个点：5V 发送端、转换器输出端和 3.3V 电源轨。正常上电状态下核对高低电平与边沿；接收端断电、发送端保持高电平时，确认 3.3V 电源没有被抬升；快速插拔或 ESD 测试后，确认转换器输出仍回到规定默认状态。仅在正常上电时看到逻辑分析仪显示 0/1，不能证明掉电和异常路径已经闭合。

\begin{longtable}{L{0.20\linewidth}L{0.28\linewidth}L{0.30\linewidth}L{0.12\linewidth}}
\toprule
方案 & 适用条件 & 必须核对的边界 & 评价 \\
\midrule
直接连接 & 接收脚明确标注 5V tolerant & 上电、掉电和所有输入模式均允许 5V & 器件明确支持时最简单 \\\tablerowrule
串联限流 & 短暂异常或手册明确允许钳位电流 & 最大钳位电流、持续时间、反灌路径 & 主要是保护，不默认作为转换 \\\tablerowrule
电阻分压 & 单向、低速、输入漏电小 & 最坏阈值、RC 边沿、功耗与掉电状态 & 简单但速度和方向受限 \\\tablerowrule
专用缓冲/转换器 & 高速、双向或电源时序复杂 & 方向控制、输出使能、掉电高阻与传播延迟 & 接口契约最完整 \\
\bottomrule
\end{longtable}

\section*{章末练习}

\begin{chapterexercise}{1}
电平判断
\exercisecheck{题目}{某输入规定 $V_{IL,max}=0.8$ V、$V_{IH,min}=2.0$ V。分别判断 0.5 V、1.4 V 和 2.7 V 属于哪一区间，并说明中间区为何不能当作第三种逻辑值。}
\end{chapterexercise}

\begin{chapterexercise}{2}
噪声余量
\exercisecheck{题目}{驱动器保证 $V_{OH,min}=2.7$ V、$V_{OL,max}=0.4$ V，接收器指标同题 1。计算高、低电平噪声余量。}
\end{chapterexercise}

\begin{chapterexercise}{3}
器件路径
\exercisecheck{题目}{用受控导通路径解释 CMOS 反相器输入为 0 和 1 时输出为何分别趋近电源与地。}
\end{chapterexercise}

\begin{chapterexercise}{4}
功耗
\exercisecheck{题目}{一个节点每秒翻转次数加倍、负载电容不变，动态功耗如何变化？若电压降为原来的 0.8 倍呢？}
\end{chapterexercise}

\begin{chapterexercise}{5}
边沿
\exercisecheck{题目}{为什么慢边沿会增加门电路短路电流和对噪声的敏感时间？施密特触发输入能改善什么？}
\end{chapterexercise}

\begin{chapterexercise}{6}
排错
\exercisecheck{题目}{板上数字输入偶发误判，按电源、阈值、边沿、串扰和测量参考地列出检查顺序。}
\end{chapterexercise}

\section*{参考解答与要点}

\begin{chapteranswer}{1}
电平判断
\answercheck{解答要点}{0.5 V 保证为低，2.7 V 保证为高，1.4 V 位于未规定区；该区的输出受器件、噪声和时间影响，不能作为稳定逻辑状态。}
\end{chapteranswer}

\begin{chapteranswer}{2}
噪声余量
\answercheck{解答要点}{高电平为 $2.7-2.0=0.7$ V，低电平为 $0.8-0.4=0.4$ V。}
\end{chapteranswer}

\begin{chapteranswer}{3}
器件路径
\answercheck{解答要点}{输入低时上拉网络导通、下拉截止，输出电容向电源充电；输入高时相反，输出经下拉网络放电到地。}
\end{chapteranswer}

\begin{chapteranswer}{4}
功耗
\answercheck{解答要点}{动态功耗近似正比于翻转率，次数加倍则加倍；电压为 0.8 倍时电压项使功耗乘 $0.8^2=0.64$。}
\end{chapteranswer}

\begin{chapteranswer}{5}
边沿
\answercheck{解答要点}{慢边沿让上下管同时部分导通更久，也延长输入停留在阈值附近的时间；施密特触发以滞回提高慢变或带噪输入的稳定性。}
\end{chapteranswer}

\begin{chapteranswer}{6}
排错
\answercheck{解答要点}{先确认电源和参考地，再量接收端实际高低电平与余量，检查上升下降时间、反射/串扰和探头接地，最后核对输入类型与阈值配置。}
\end{chapteranswer}
